\documentclass[UTF8, 12pt, fontset=fandol, oneside]{ctexart}
\usepackage{researchnotes}

\title{凸优化：逼近与拟合}
\author{}
\date{}

\begin{document}

\maketitle

\section*{6\quad 逼近与拟合}

\subsection*{6.1\quad 范数逼近}

\subsection*{6.1.1\quad 基本的范数逼近问题}

最简单的范数逼近问题是具有下列形式的无约束问题
\begin{equation}
\begin{aligned}
\text{minimize}\quad & \|Ax-b\|,
\end{aligned}
\tag{6.1}
\end{equation}
其中 $A\in\mathbb{R}^{m\times n}$ 和 $b\in\mathbb{R}^m$ 是问题的数据，$x\in\mathbb{R}^n$ 是变量，而 $\|\cdot\|$ 是 $\mathbb{R}^m$ 上的一种范数。范数逼近问题的解有时又被称为 $Ax\approx b$ 在范数 $\|\cdot\|$ 下的近似解。向量
\[
r=Ax-b
\]
称为这个问题的残差；其分量有时也称为关于 $x$ 的个体残差。

范数逼近问题 (6.1) 是一个可解的凸问题，也就是说，总是存在至少一个最优解。当且仅当 $b\in\mathcal{R}(A)$ 时，其最优值为零；但是，更有趣也更有用的是 $b\notin\mathcal{R}(A)$ 的情况。我们可以不失一般性地假设 $A$ 的列向量独立；特别地，设 $m\geq n$。当 $m=n$ 时，最优解 $A^{-1}b$ 可以简单地得到，因此我们可以假设 $m>n$。

\noindent\textbf{逼近的解释}\quad
通过将 $Ax$ 表示为
\[
Ax=x_1a_1+\cdots+x_na_n,
\]
其中 $a_1,\cdots,a_n\in\mathbb{R}^m$ 为 $A$ 的列，我们可以看出，范数逼近问题的目标是用 $A$ 的列的线性组合，尽可能准确地逼近或拟合向量 $b$，其偏差由范数 $\|\cdot\|$ 度量。

这一逼近问题也称为回归问题。在此背景下，向量 $a_1,\cdots,a_n$ 称为回归量，设 $x$ 是问题的一个最优解，称向量 $x_1a_1+\cdots+x_na_n$ 为 $b$（向回归量）的回归。

\noindent\textbf{估计的解释}\quad
一个与范数逼近问题密切相关的解释出现在基于不完全线性向量测量值进行的参数估计问题中。考虑线性测量模型
\[
y=Ax+v,
\]
其中 $y\in\mathbb{R}^m$ 为测量值向量，$x\in\mathbb{R}^n$ 为待估计的参数向量，$v\in\mathbb{R}^m$ 为未知的测量误差（假设在范数 $\|\cdot\|$ 度量下很小）。相应的估计问题是在给定 $y$ 的情况下对 $x$ 是什么进行合理的猜测。

如果我们猜测 $x$ 的值为 $\hat{x}$，也就是隐含地猜测 $v$ 的值为 $y-A\hat{x}$。假设 $v$（在 $\|\cdot\|$ 度量下）越小越可信，那么对于 $x$ 的最有可信的猜测是
\[
\hat{x}=\arg\min_z \|Az-y\|.
\]
（这些想法可以在统计学的框架中进行更加正规的表述，参见第 7 章。）

\noindent\textbf{几何的解释}\quad
考虑子空间 $\mathcal{A}=\mathcal{R}(A)\subseteq\mathbb{R}^m$ 和一个点 $b\in\mathbb{R}^m$。点 $b$ 向子空间 $\mathcal{A}$ 的投影是 $\mathcal{A}$ 中在 $\|\cdot\|$ 下最靠近 $b$ 的点，也就是说，它是下列问题的任意最优解
\[
\begin{aligned}
\text{minimize}\quad & \|u-b\|\\
\text{subject to}\quad & u\in\mathcal{A}.
\end{aligned}
\]
将 $\mathcal{R}(A)$ 中的元素参数化为 $u=Ax$，我们可以看出求解范数逼近问题 (6.1) 等价于计算 $b$ 向 $\mathcal{A}$ 的投影。

\noindent\textbf{设计的解释}\quad
我们可以将范数逼近问题 (6.1) 解释为一个最优设计问题。$x_1,\cdots,x_n$ 为 $n$ 个设计变量，其数值待定。向量 $y=Ax$ 给出了表示 $m$ 个结果的向量，我们假设它是设计变量 $x$ 的线性函数。向量 $b$ 为目标向量或期望结果。目标是选择一个设计向量尽可能地接近期望的结果，也就是说 $Ax\approx b$。我们可以将残差向量 $r$ 解释为实际结果（即 $Ax$）与期望或目标（即 $b$）之间的偏差。如果我们用真实结果与期望目标之间偏差的范数来度量设计的质量，那么参数逼近问题 (6.1) 也就是寻找最佳设计的问题。

\noindent\textbf{加权范数逼近问题}\quad
范数逼近问题的一个扩展是加权范数逼近问题
\[
\begin{aligned}
\text{minimize}\quad & \|W(Ax-b)\|,
\end{aligned}
\]
其问题数据 $W\in\mathbb{R}^{m\times m}$ 称为权矩阵。权矩阵通常是对角的，在这种情况下，它给出了对残差向量 $r=Ax-b$ 分量之间不同的相对强调程度。

加权范数问题可以看做关于范数 $\|\cdot\|$ 和数据 $\tilde{A}=WA,\tilde{b}=Wb$ 的范数逼近问题，因此，仍然可以被视为标准范数逼近问题 (6.1)。另一方面，加权范数逼近问题也可以被看做关于数据 $A$ 和 $b$ 的范数估计问题，其范数为 $W$-加权范数，定义如下
\[
\|z\|_W=\|Wz\|
\]
（这里假设 $W$ 非奇异）。

\noindent\textbf{最小二乘逼近}\quad
最常见的范数逼近问题采用 Euclid 或 $\ell_2$-范数。通过将目标函数平方，我们得到等价问题，称为最小二乘逼近问题，
\[
\begin{aligned}
\text{minimize}\quad & \|Ax-b\|_2^2=r_1^2+r_2^2+\cdots+r_m^2,
\end{aligned}
\]
其目标函数为残差平方和。通过将目标函数表示为凸二次函数，
\[
f(x)=x^TA^TAx-2b^TAx+b^Tb.
\]
这个问题可以解析地求解。点 $x$ 极小化 $f$，当且仅当
\[
\nabla f(x)=2A^TAx-2A^Tb=0,
\]
即 $x$ 满足
\[
A^TAx=A^Tb,
\]
这个方程称为正规方程，并且总是有解的。因为我们假设 $A$ 的列向量是独立的，所以最小二乘逼近问题总有唯一解 $x=(A^TA)^{-1}A^Tb$。

\noindent\textbf{Chebyshev 或极小极大逼近}\quad
当使用 $\ell_\infty$-范数时，范数逼近问题
\[
\begin{aligned}
\text{minimize}\quad & \|Ax-b\|_\infty=\max\{|r_1|,\cdots,|r_m|\}
\end{aligned}
\]
称为 Chebyshev 逼近问题。因为我们试图极小化最大（绝对值）残差，因此又称为极小极大逼近问题。Chebyshev 逼近问题可以描述为线性规划
\[
\begin{aligned}
\text{minimize}\quad & t\\
\text{subject to}\quad & -t\mathbf{1}\preceq Ax-b\preceq t\mathbf{1},
\end{aligned}
\]
其变量为 $x\in\mathbb{R}^n$ 和 $t\in\mathbb{R}$。

\noindent\textbf{残差绝对值之和逼近}\quad
如果采用 $\ell_1$-范数，范数逼近问题
\[
\begin{aligned}
\text{minimize}\quad & \|Ax-b\|_1=|r_1|+\cdots+|r_m|
\end{aligned}
\]
称为残差（绝对值）和逼近问题，或者，在估计领域，称为一种鲁棒估计器（其原因马上会明了）。如同 Chebyshev 逼近问题一样，$\ell_1$-范数逼近问题可以描述为线性规划
\[
\begin{aligned}
\text{minimize}\quad & \mathbf{1}^Tt\\
\text{subject to}\quad & -t\preceq Ax-b\preceq t,
\end{aligned}
\]
其变量为 $x\in\mathbb{R}^n$ 和 $t\in\mathbb{R}^m$。

\subsection*{6.1.2\quad 罚函数逼近}

对于 $1\leq p<\infty$，$\ell_p$-范数逼近问题的目标函数为
\[
\left(|r_1|^p+\cdots+|r_m|^p\right)^{1/p}.
\]
如同最小二乘问题一样，我们可以考虑目标函数为
\[
|r_1|^p+\cdots+|r_m|^p,
\]
的等价问题，这一目标函数是残差的可分、对称函数。特别地，目标函数值仅取决于残差的幅值分布，即排序的残差。

我们将考虑 $\ell_p$-范数逼近问题的一个有用推广，其目标函数仅仅取决于残差的幅值分布。这个罚函数逼近问题具有形式
\begin{equation}
\begin{aligned}
\text{minimize}\quad & \phi(r_1)+\cdots+\phi(r_m)\\
\text{subject to}\quad & r=Ax-b,
\end{aligned}
\tag{6.2}
\end{equation}
其中 $\phi:\mathbb{R}\to\mathbb{R}$ 称为（残差）罚函数。设 $\phi$ 为凸函数，则罚函数逼近问题是一个凸优化问题。在很多情况下，罚函数 $\phi$ 是对称、非负的，并且满足 $\phi(0)=0$。但是，在我们的分析中，并不需要这些性质。

\noindent\textbf{解释}\quad
我们可以将罚函数逼近问题 (6.2) 解释如下。对于 $x$ 的选择，我们得到了 $b$ 的一个逼近 $Ax$，也得到了相应的残差向量 $r$。罚函数通过 $\phi(r_i)$ 评价残差每一个分量的费用或惩罚；总体惩罚就是每个残差的罚函数之和，即 $\phi(r_1)+\cdots+\phi(r_m)$。$x$ 的不同选择导致不同的残差，因此有不同的总体惩罚。在罚函数逼近问题中，我们极小化由残差带来的总体惩罚。

\noindent\textbf{例 6.1}\quad
一些常见的罚函数及其逼近问题。

采用 $\phi(u)=|u|^p$，其中 $p\geq 1$，罚函数逼近问题等价于 $\ell_p$-范数逼近问题。特别地，二次罚函数 $\phi(u)=u^2$ 导致最小二乘或 Euclid 范数逼近，绝对值罚函数 $\phi(u)=|u|$ 导致 $\ell_1$-范数逼近。

带有死区的线性罚函数（死区宽度为 $a>0$）由下式给出
\[
\phi(u)=
\begin{cases}
0 & |u|\leq a\\
|u|-a & |u|>a.
\end{cases}
\]
带有死区的线性函数对于小于 $a$ 的残差不进行惩罚。

对数障碍罚函数（极限为 $a>0$）具有下面的形式，
\[
\phi(u)=
\begin{cases}
-a^2\log(1-(u/a)^2) & |u|<a\\
\infty & |u|\geq a.
\end{cases}
\]
对数障碍罚函数对于大于 $a$ 的残差给予无穷的惩罚。

图 6.1 给出了带有死区的线性罚函数、对数障碍罚函数及二次罚函数的图像。需要注意的是，当 $|u/a|\leq 0.25$ 时，对数罚函数与二次罚函数非常地接近（参见习题 6.1）。

将罚函数缩放某个正数倍不会影响罚函数逼近问题的解，因为这仅仅对目标函数进行了伸缩变换。但是罚函数的形状对罚函数逼近问题的解有很大的影响。粗略地，$\phi(u)$ 度量了我们不喜欢残差值 $u$ 的程度。如果对于比较小的 $u$，$\phi$ 值很小，表示我们不大（或完全不）关心具有这些值的残差。如果当 $u$ 变大时，$\phi(u)$ 增长迅速，表示我们对于大的误差非常厌恶；如果在某些区间外，$\phi$ 变成无穷，表示在这些区间外的残差是无法接受的。这些简单的解释给出了对罚函数逼近问题的解的深入理解，同时对如何选择罚函数也有所启示。

作为一个例子，我们比较 $\ell_1$-范数和 $\ell_2$-范数逼近，也就是分别对应 $\phi_1(u)=|u|$ 和 $\phi_2(u)=u^2$ 的罚函数逼近问题。对于 $|u|=1$，这两个函数给予同样的罚。对于小的 $u$，我们有 $\phi_1(u)\gg \phi_2(u)$，因此，相对于 $\ell_2$-范数逼近，$\ell_1$-范数逼近对较小的残差给予较大的重视。对于大的 $u$，我们有 $\phi_2(u)\gg \phi_1(u)$，因此，相比 $\ell_2$-范数逼近，$\ell_1$-范数逼近给予大的残差的权重较小。对于大小残差的相对权重的不同，反应在相关逼近问题的解上。相对于 $\ell_2$-范数，$\ell_1$-范数逼近问题最优残差的幅值分布趋向于有更多零或非常小的残差。相反，$\ell_2$-范数的解会趋向于有较少大的残差（因为相比于 $\ell_1$-范数，$\ell_2$-范数逼近问题中，大的残差将带来更大的惩罚）。

\noindent\textbf{例子}\quad
下面的例子将说明上述讨论。我们选取矩阵 $A\in\mathbb{R}^{100\times 30}$ 和向量 $b\in\mathbb{R}^{100}$（数值随机产生，但其结果是典型的），然后比较 $Ax\approx b$ 的 $\ell_1$-范数、$\ell_2$-范数逼近结果，我们还将比较带有死区的线性罚函数（$a=0.5$）以及对数障碍罚函数（$a=1$）的结果。图 6.2 显示了这四种罚函数及相应最优残差的分布。从罚函数的图像中，我们注意到：

$\ell_1$-范数罚函数对小的残差给予最重的权，而对大的残差给予最小的权。

$\ell_2$-范数罚函数对小的残差给予非常小的权，而对大的残差给予很重的权。

带有死区的线性罚函数对于小于 $0.5$ 的残差不予惩罚，而对于较大的残差，其惩罚也相对较小。

对数障碍罚函数对于小的残差，给予非常近似于 $\ell_2$-范数罚函数的惩罚，而对于大于大约 $0.8$ 的残差给予非常重的惩罚，并对大于 1 的残差给予无限的权。

从幅值分布可以看出一些性质：

对于 $\ell_1$-最优解，很多残差为零或非常小。$\ell_1$-最优解也含有很多相对较大的残差。

$\ell_2$-范数逼近结果有很多适度的残差，而较大的残差相对较少。

对于带有死区的线性罚函数，我们可以看出有很多等于 $\pm 0.5$ 的残差，正好处于``自由''区域的边缘，因为对这些残差并未施加以任何惩罚。

对于对数障碍罚函数，我们可以看出，没有任何幅值大于 1 的残差，除此以外，其残差分布与 $\ell_2$-范数逼近的结果很相似。

\noindent\textbf{对于野值或大误差的灵敏性}\quad
在估计或回归领域，野值是指具有相对很大的噪声 $v_i$ 的测量值 $y_i=a_i^Tx+v_i$。在与故障数据或错误测量相关的领域中，这种情况很普遍。当野值存在时，$x$ 的任意估计结果都会产生含有较大分量的残差向量。理想地，我们希望能够猜出哪些测量值是野值，然后或者把它们从估计过程中移除，或者在估计时大大降低它们的权。（但是，我们不能对大的残差给予零惩罚，这是因为如果这样做了，那么最优点一定会趋向于使得所有残差变大以取得总惩罚为零的结果。）这一想法可以通过罚函数逼近来实现，比如设计这样的罚函数，
\begin{equation}
\phi(u)=
\begin{cases}
u^2 & |u|\leq M\\
M^2 & |u|>M,
\end{cases}
\tag{6.3}
\end{equation}
如图 6.3 所示。对于任意小于 $M$ 的残差，这样的罚函数与最小二乘相吻合，但是对于大于 $M$ 的残差，无论大于多少，它都给予一个固定的权。换言之，大于 $M$ 的残差被忽略了；它们被假设与野值或不好的数据相关联。不幸的是，罚函数 (6.3) 不是凸的，相应的罚函数逼近问题将变成一个很难求解的组合优化问题。

基于估计方法的罚函数，其对野值的灵敏度取决于罚函数对于大的残差的（相对）惩罚值。如果限定为凸的罚函数（将得到凸优化问题），最为不灵敏的罚函数是那些 $\phi(u)$ 线性增长的函数，即对于较大的 $u$，$\phi(u)$ 近似于 $|u|$。有这样性质的罚函数有时也称为鲁棒的，因为相应的罚函数逼近方法对于野值或大的误差，相比于例如最小二乘这样的方法，要不灵敏得多。

鲁棒罚函数的一个显然的例子是 $\phi(u)=|u|$，对应于 $\ell_1$-范数逼近。另一个例子是鲁棒最小二乘或 Huber 罚函数，
\begin{equation}
\phi_{\mathrm{hub}}(u)=
\begin{cases}
u^2 & |u|\leq M\\
M(2|u|-M) & |u|>M,
\end{cases}
\tag{6.4}
\end{equation}
其图像见图 6.4。当残差小于 $M$ 时，这个罚函数与最小二乘的罚函数相同，而对于大的残差，它又恢复为类似 $\ell_1$ 的线性增长。Huber 罚函数也可以被视为是对野值罚函数 (6.3) 的凸近似：无论 $|u|\leq M$ 还是 $|u|>M$，它们都很相似，即 Huber 罚函数是最接近野值罚函数 (6.3) 的凸函数。

\noindent\textbf{例 6.2}\quad
鲁棒回归。图 6.5 显示了平面上的 42 个点 $(t_i,y_i)$，其中有两个明显的野值（一个在左上，一个在右下）。虚线表明了用直线 $f(t)=\alpha+\beta t$ 对这些点的最小二乘逼近结果。系数 $\alpha$ 和 $\beta$ 是通过求解最小二乘问题
\[
\begin{aligned}
\text{minimize}\quad & \sum_{i=1}^{42}(y_i-\alpha-\beta t_i)^2
\end{aligned}
\]
得到的，这个问题的变量为 $\alpha$ 和 $\beta$。显然，最小二乘逼近的结果偏离了这些点，向着两个野值进行了偏转。

实线显示了鲁棒最小二乘逼近的结果，通过极小化 Huber 罚函数得到，即
\[
\begin{aligned}
\text{minimize}\quad & \sum_{i=1}^{42}\phi_{\mathrm{hub}}(y_i-\alpha-\beta t_i),
\end{aligned}
\]
其中 $M=1$。这样的逼近受野值的影响小了很多。

因为 $\ell_1$-范数逼近是（凸）罚函数逼近方法中对野值最为鲁棒的方法，$\ell_1$-范数逼近也称为鲁棒估计或鲁棒回归。$\ell_1$-范数估计的鲁棒性也可以在统计框架下理解，参见第 339 页。

\noindent\textbf{小残差与 $\ell_1$-范数逼近}\quad
我们也关注小的残差。最小二乘逼近对小的残差给予很小的权，因为当 $u$ 很小时，$\phi(u)=u^2$ 非常小。类似带有死区的线性罚函数，一些罚函数对最小残差给予零权。对于那些对小的残差给予非常小的惩罚的罚函数，我们期望它们能够使最优残差比较小，但不要求非常小。粗略地，没有或很少有动力去驱动小残差变得更小。

相反，对较小残差给予相对较大权的罚函数，例如 $\phi(u)=|u|$，即相应的 $\ell_1$-范数逼近，往往得到含有很多非常小、甚至恰好为零的分量的最优残差。这意味着，在 $\ell_1$-范数逼近中，通常有很多等式被精确地满足，也就是说，对于很多 $i$ 有 $a_i^Tx=b_i$。这一现象可从图 6.2 中看出。

\subsection*{6.1.3\quad 带有约束的逼近}

可以在基本的范数逼近问题 (6.1) 中添加约束。当这些约束为凸约束时，得到的问题是凸的。有多种原因产生这些约束。

在逼近问题中，约束可以用来排除对向量 $b$ 的特定的、不可接受的逼近，或者确保 $Ax$ 的逼近结果满足特定的性质。

在逼近问题中，约束可以作为待估计向量 $x$ 的先验知识，或者来源于估计误差 $v$ 的先验知识。

确定点 $b$ 向比子空间更复杂的集合（例如，锥或多面体）的投影时，约束具有相应的几何意义。

一些例子将说明上述讨论。

\noindent\textbf{变量的非负约束}\quad
我们可以在基本的范数逼近问题中添加约束 $x\succeq 0$：
\[
\begin{aligned}
\text{minimize}\quad & \|Ax-b\|\\
\text{subject to}\quad & x\succeq 0.
\end{aligned}
\]
当我们估计已知非负的参数向量 $x$ 时，例如，能量、强度或比率，需要在估计中添加非负约束。其几何意义是计算向量 $b$ 向由 $A$ 的列向量组成的锥的投影。我们也可以将这一问题解释为用 $A$ 的列向量的非负组合（即锥组合）去逼近 $b$。

\noindent\textbf{变量范围}\quad
这里我们添加约束 $l\preceq x\preceq u$，其中 $l,u\in\mathbb{R}^n$ 为问题的参数：
\[
\begin{aligned}
\text{minimize}\quad & \|Ax-b\|\\
\text{subject to}\quad & l\preceq x\preceq u.
\end{aligned}
\]
在估计问题中，变量范围作为先验知识提出，表明了每个变量所处的区间。其几何解释是我们在确定 $b$ 向集合 $\{Ax\mid l\preceq x\preceq u\}$ 的投影。

\noindent\textbf{概率分布}\quad
我们可以添加约束使得 $x$ 满足 $x\succeq 0,\mathbf{1}^Tx=1$：
\[
\begin{aligned}
\text{minimize}\quad & \|Ax-b\|\\
\text{subject to}\quad & x\succeq 0,\quad \mathbf{1}^Tx=1.
\end{aligned}
\]
这常见于对比例或相对频率（它们非负且和为 1）的估计中。这也可以解释为用 $A$ 的列向量的凸组合来逼近 $b$。（对于概率估计，我们将在 $\S 7.2$ 中进行更多的讨论。）

\noindent\textbf{范数球约束}\quad
在基本的范数逼近问题中，我们可以添加约束使得 $x$ 位于一个范数球中：
\[
\begin{aligned}
\text{minimize}\quad & \|Ax-b\|\\
\text{subject to}\quad & \|x-x_0\|\leq d,
\end{aligned}
\]
其中 $x_0$ 和 $d$ 为问题的参数。这一约束可以由于下述原因而添加：

在估计问题中，$x_0$ 是对参数 $x$ 的一个先验的猜测，而 $d$ 是估计值距离先验猜测的最大可能误差。我们对于参数 $x$ 的估计是在所有可能的候选值（即满足 $\|z-x_0\|\leq d$ 的 $z$）中选择与测量数据匹配最好的 $\hat{x}$（即极小化 $\|Az-b\|$）。

约束 $\|x-x_0\|\leq d$ 可以表示信赖域。线性关系 $y=Ax$ 仅仅是某些非线性关系 $y=f(x)$ 的近似，这种近似在 $x_0$ 的附近，具体地，$\|x-x_0\|\leq d$ 成立。于是问题变成在这些 $x$ 上，即在模型 $y=Ax$ 可以被信赖的区域中，极小化 $\|Ax-b\|$。

这些想法在正则化的内容中也会被提及，参见 $\S 6.3.2$。

\section*{6.2\quad 最小范数问题}

基本的最小范数问题具有下列形式
\begin{equation}
\begin{aligned}
\text{minimize}\quad & \|x\|\\
\text{subject to}\quad & Ax=b,
\end{aligned}
\tag{6.5}
\end{equation}
其中数据为 $A\in\mathbb{R}^{m\times n}$ 和 $b\in\mathbb{R}^m$，变量为 $x\in\mathbb{R}^n$，$\|\cdot\|$ 为 $\mathbb{R}^n$ 上的一种范数。这个问题的解称为 $Ax=b$ 的最小范数解，如果 $Ax=b$ 有解，这样的解总是存在的。当然，最小范数问题是一个凸优化问题。

我们可以不失一般性地假设 $A$ 的行独立，因此 $m\leq n$。当 $m=n$ 时，唯一的可行解是 $x=A^{-1}b$；只有当 $m<n$，即方程 $Ax=b$ 不定时，最小范数问题才有意义。

\noindent\textbf{重构为范数逼近问题}\quad
通过消去等式约束，可以将最小范数问题 (6.5) 表示为一个范数逼近问题。令 $x_0$ 为 $Ax=b$ 的一个任意解，令 $Z\in\mathbb{R}^{n\times k}$ 的列为 $A$ 的零空间的基。于是，$Ax=b$ 的通解可以表示为 $x_0+Zu$，其中 $u\in\mathbb{R}^k$。最小范数问题 (6.5) 可以写为
\[
\begin{aligned}
\text{minimize}\quad & \|x_0+Zu\|,
\end{aligned}
\]
其变量为 $u\in\mathbb{R}^k$，这是一个范数逼近问题。特别地，（如果演绎正确，）关于范数逼近问题的讨论和分析都可以应用于最小范数问题。

\noindent\textbf{控制或设计的解释}\quad
我们可以将最小范数问题 (6.5) 解释为最优设计或最优控制问题。$x_1,\cdots,x_n$ 为 $n$ 个设计变量，其值待定。在控制领域，变量 $x_1,\cdots,x_n$ 表示输入，其值待我们选定。向量 $y=Ax$ 给出了设计 $x$ 对应的 $m$ 个属性或结果，我们假设它们是设计变量 $x$ 的线性函数。$m<n$ 个方程 $Ax=b$ 表示了设计的 $m$ 个规范或要求。因为 $m<n$，这种设计是不定的，具有 $n-m$ 个自由度（假设 $A$ 的秩为 $m$）。

在所有满足规格的设计中，最小范数问题选择了范数 $\|\cdot\|$ 度量下的最小方案。这可以认为是最为有效的设计，因为它用最小的可能值 $x$ 达到了要求 $Ax=b$。

\noindent\textbf{估计解释}\quad
假设 $x$ 为待估计的参数向量。我们有 $m<n$ 个很好的（没有噪声的）线性测量值，由 $Ax=b$ 给出。因为测量值少于待估计的参数，我们的测量值不足以完全确定 $x$。任意满足 $Ax=b$ 的参数向量 $x$ 都符合我们的测量值。

为了不进行更多的测量而得到 $x$ 的一个好的估计，我们必须用到先验信息。假设我们有先验信息，或者单纯假设，$x$ 很可能较小（在 $\|\cdot\|$ 度量下）。在所有满足测量值 $Ax=b$ 的参数向量中，最小范数问题选择了最小的（因此，也是最有可能的）向量作为参数向量 $x$ 的估计。（最小范数问题的统计解释可见第 345 页。）

\noindent\textbf{几何解释}\quad
我们也可以对最小二乘问题 (6.5) 给出简单的几何解释。可行集 $\{x\mid Ax=b\}$ 是仿射的，目标函数是 $x$ 和 0（在范数 $\|\cdot\|$ 度量下）的距离。最小范数问题在仿射集合中寻找距离 0 最近的点，即寻找 0 向仿射集合 $\{x\mid Ax=b\}$ 的投影。

\noindent\textbf{线性方程组的最小二乘解}\quad
最常见的最小范数问题采用 Euclid 或 $\ell_2$-范数。通过将目标函数平方，我们得到等价问题
\[
\begin{aligned}
\text{minimize}\quad & \|x\|_2^2\\
\text{subject to}\quad & Ax=b,
\end{aligned}
\]
其唯一解称为方程组 $Ax=b$ 的最小二乘解。类似最小二乘逼近问题，这一问题也可以被解析地求解。引入对偶变量 $\nu\in\mathbb{R}^m$，最优性条件为
\[
2x^\star+A^T\nu^\star=0,\qquad Ax^\star=b,
\]
这是一组线性方程并容易求解。通过第一个方程，我们可以得到 $x^\star=-(1/2)A^T\nu^\star$；将其带入第二个方程，我们有 $-(1/2)AA^T\nu^\star=b$ 并且得到结论
\[
\nu^\star=-2(AA^T)^{-1}b,\qquad x^\star=A^T(AA^T)^{-1}b.
\]
（因为 $\mathbf{rank}\,A=m<n$，矩阵 $AA^T$ 可逆。）

\noindent\textbf{最小罚问题}\quad
最小范数问题 (6.5) 一个有用的变形是最小罚问题
\begin{equation}
\begin{aligned}
\text{minimize}\quad & \phi(x_1)+\cdots+\phi(x_n)\\
\text{subject to}\quad & Ax=b,
\end{aligned}
\tag{6.6}
\end{equation}
其中 $\phi:\mathbb{R}\to\mathbb{R}$ 是凸和非负的并且满足 $\phi(0)=0$。罚函数的值 $\phi(u)$ 量化了我们对 $x$ 的某一分量具有值 $u$ 的厌恶程度；于是，在约束 $Ax=b$ 下，最小罚问题找到了具有最小总惩罚的 $x$。

通过将（罚函数逼近问题中的）残差 $r$ 的幅值分布替换为（最小罚问题中的）$x$ 的幅值分布，罚函数逼近问题中所有关于罚函数的讨论和解释都可以转换到最小罚问题上。

\noindent\textbf{通过最小 $\ell_1$-范数得到稀疏解}\quad
重新考虑第 293 页的讨论，$\ell_1$-范数逼近对小的残差给予相对较大的权，因此，将导致很多最优残差很小，或者甚至是零。类似的效果也发生在最小范数领域。最小 $\ell_1$-范数问题
\[
\begin{aligned}
\text{minimize}\quad & \|x\|_1\\
\text{subject to}\quad & Ax=b
\end{aligned}
\]
得到的解 $x$ 趋向于有很多等于零的分量。换言之，最小 $\ell_1$-范数问题趋向于得到 $Ax=b$ 的稀疏解，常常有 $m$ 个非零分量。

容易找到 $Ax=b$ 的只有 $m$ 个非零分量的解。（从 $1,\cdots,n$ 中）任意选择 $m$ 个指标作为 $x$ 的非零分量的下标。方程 $Ax=b$ 退化为 $\tilde{A}\tilde{x}=b$，其中 $\tilde{A}$ 为选定的 $A$ 的列组成的 $m\times m$ 子矩阵，而 $\tilde{x}\in\mathbb{R}^m$ 为 $x$ 的子向量，包含了 $m$ 个选定的分量。如果 $\tilde{A}$ 非奇异，那么我们可以得到 $\tilde{x}=\tilde{A}^{-1}b$，从而给出一个有 $m$ 或更少非零分量的可行解。如果 $\tilde{A}$ 奇异并且 $b\notin\mathcal{R}(\tilde{A})$，那么方程 $\tilde{A}\tilde{x}=b$ 不可解，这意味着对于选定的非零分量，不存在可行解。如果 $\tilde{A}$ 奇异且 $b\in\mathcal{R}(\tilde{A})$，那么存在可行解，其非零分量个数少于 $m$。

这一方法可以用来寻找有 $m$ 个（或更少）非零元素的最小的 $x$，但是，这通常需要考虑并比较所有 $n!/(m!(n-m)!)$ 种组合：从 $x$ 的 $n$ 个系数中选择 $m$ 个非零系数。另一方面，求解最小 $\ell_1$-范数问题，给出了好的启发式算法用以寻找 $Ax=b$ 的稀疏的、较小的解。

\section*{6.3\quad 正则化逼近}

\subsection*{6.3.1\quad 双准则式}

在正则化逼近的基本形式中，我们的目标是寻找向量 $x$ 使其较小（如果可能的话），同时使得残差 $Ax-b$ 小。自然地，这可以描述为双目标的（凸）向量优化问题，这两个目标是 $\|Ax-b\|$ 和 $\|x\|$：
\begin{equation}
\begin{aligned}
\text{minimize（关于 $\mathbb{R}_+^2$）}\quad & (\|Ax-b\|,\|x\|).
\end{aligned}
\tag{6.7}
\end{equation}
这两个范数可能是不同的：第一个在 $\mathbb{R}^m$ 中，用以度量残差的规模；第二个在 $\mathbb{R}^n$ 中，用以度量 $x$ 的规模。

可以通过多种方法找到这两个目标之间的最优权衡，然后给出 $\|Ax-b\|$ 关于 $\|x\|$ 的最优权衡曲线，显示了为使一个目标减小而另一个目标必须增加的量。容易表述 $\|Ax-b\|$ 和 $\|x\|$ 之间最优权衡曲线的一个端点：$\|x\|$ 的最小值为零，只有当 $x=0$ 时达到。对于 $x$ 的这个值，残差范数的值为 $\|b\|$。

权衡曲线另一个端点的描述要复杂一些。用 $C$ 表示 $\|Ax-b\|$ 的最小解集合（没有关于 $\|x\|$ 的约束）。那么 $C$ 中的任意最小范数点都是 Pareto 最优的，对应权衡曲线上的另一个端点。换言之，在这个端点处的 Pareto 最优点由范数极小化 $\|Ax-b\|$ 的最小解给出。如果是 Euclid 范数，这个 Pareto 最优解是唯一的，由 $x=A^\dagger b$ 给出，其中 $A^\dagger$ 为 $A$ 的伪逆。（参见 $\S 4.7.6$，第 176 页以及 $\S A.5.4$。）

\subsection*{6.3.2\quad 正则化}

正则化是求解双准则问题 (6.7) 的一个常用的标量化方法。正则化的一种形式是极小化目标函数的加权和：
\begin{equation}
\begin{aligned}
\text{minimize}\quad & \|Ax-b\|+\gamma\|x\|,
\end{aligned}
\tag{6.8}
\end{equation}
其中 $\gamma>0$ 为问题参数。当 $\gamma$ 在 $(0,\infty)$ 上变化时，(6.8) 的解遍历了最优权衡曲线。

正则化的另一个常用方法是极小化加权范数平方和，特别是在使用 Euclid 范数的情形下。这一方法是对变化的 $\delta>0$ 求解
\begin{equation}
\begin{aligned}
\text{minimize}\quad & \|Ax-b\|^2+\delta\|x\|^2.
\end{aligned}
\tag{6.9}
\end{equation}
这些正则化的逼近问题都可以用来求解双目标问题，通过增加与 $x$ 的范数相关的附加项或惩罚来使得 $\|Ax-b\|$ 和 $\|x\|$ 均很小。

\noindent\textbf{解释}\quad
正则化在很多领域都得到了应用。在估计领域，附加的对大的 $\|x\|$ 的惩罚可以解释为我们的先验知识：$\|x\|$ 不是很大。在最优设计领域，附加项将使用较大设计变量的成本添加到了偏离目标的成本中。

$\|x\|$ 应当较小这个约束也反映在建模问题中。例如，$y=Ax$ 可能仅仅是 $x$ 和 $y$ 之间真实关系 $y=f(x)$ 的一个很好的逼近。为使得 $f(x)\approx b$，我们希望 $Ax\approx b$ 同时需要 $x$ 较小以保证 $f(x)\approx Ax$。

在 $\S 6.4.1$ 和 $\S 6.4.2$ 中，我们将看到，正则化可以用来考虑矩阵 $A$ 的变化。粗略地，大的 $x$ 会使得 $A$ 上的变动引起 $Ax$ 的大的变动，因此，需要避免。

正则化也可以用于在矩阵 $A$ 稀疏的条件下求解线性方程 $Ax=b$。在 $A$ 的条件数较坏，或者甚至是奇异时，正则化在求解方程组（即，使得 $\|Ax-b\|$ 为零）和保持 $x$ 在合理的规模间取得折中。

正则化也可归结到统计的框架中，参见 $\S 7.1.2$。

\noindent\textbf{Tikhonov 正则化}\quad
最常用的正则化方法基于式 (6.9) 并利用 Euclid 范数，这将得到一个（凸）二次优化问题：
\begin{equation}
\begin{aligned}
\text{minimize}\quad & \|Ax-b\|_2^2+\delta\|x\|_2^2
=x^T(A^TA+\delta I)x-2b^TAx+b^Tb.
\end{aligned}
\tag{6.10}
\end{equation}
这个 Tikhonov 正则化问题有解析解
\[
x=(A^TA+\delta I)^{-1}A^Tb.
\]
因为对于任意 $\delta>0$，都有 $A^TA+\delta I\succ 0$，所以 Tikhonov 正则化的最小二乘解不需要对矩阵 $A$ 的秩（或维数）做出假设。

\noindent\textbf{光滑正则化}\quad
正则化的想法，即向目标函数添加惩罚大的 $x$ 的项，可以有很多扩展。其中一个有用的扩展是用 $\|Dx\|$ 代替 $\|x\|$ 作为正则化的项，以替代 $\|x\|$。在很多应用中，矩阵 $D$ 代表近似的微分或二阶微分算子，因此 $\|Dx\|$ 代表 $x$ 的变化的度量或者其光滑度。

例如，设向量 $x\in\mathbb{R}^n$ 表示一些连续物理参数的值，例如，温度。设这些值分布在区间 $[0,1]$ 上：$x_i$ 为点 $i/n$ 处的温度。在 $i/n$ 附近梯度或一阶导数的简单近似由 $n(x_{i+1}-x_i)$ 给出，而二阶导数的简单近似由二阶差分
\[
n\left(n(x_{i+1}-x_i)-n(x_i-x_{i-1})\right)=n^2(x_{i+1}-2x_i+x_{i-1})
\]
给出。如果 $\Delta$ 是（三对角、Toeplitz）矩阵
\[
\Delta=n^2
\begin{bmatrix}
1 & -2 & 1 & 0 & \cdots & 0 & 0 & 0 & 0\\
0 & 1 & -2 & 1 & \cdots & 0 & 0 & 0 & 0\\
0 & 0 & 1 & -2 & \cdots & 0 & 0 & 0 & 0\\
\vdots & \vdots & \vdots & \vdots & & \vdots & \vdots & \vdots & \vdots\\
0 & 0 & 0 & 0 & \cdots & -2 & 1 & 0 & 0\\
0 & 0 & 0 & 0 & \cdots & 1 & -2 & 1 & 0\\
0 & 0 & 0 & 0 & \cdots & 0 & 1 & -2 & 1
\end{bmatrix}
\in\mathbb{R}^{(n-2)\times n},
\]
那么，$\Delta x$ 表示了对参数二阶导数的近似，因此，$\|\Delta x\|_2^2$ 代表对参数在 $[0,1]$ 区间上平均平方曲率的度量。

Tikhonov 正则化问题
\[
\begin{aligned}
\text{minimize}\quad & \|Ax-b\|_2^2+\delta\|\Delta x\|_2^2
\end{aligned}
\]
可以用来在目标 $\|Ax-b\|^2$ 和 $\|\Delta x\|^2$ 间进行权衡。前者可以看成是拟合的度量，或与实验数据的一致程度；后者则（近似地）代表相关物理量的平均平方曲率。参数 $\delta$ 用以控制所需要的正则化的量，或者，给出拟合-光滑度的最优权衡曲线。

我们也可以添加多个正则化项，比如，与光滑程度和规模相关的项，例如
\[
\begin{aligned}
\text{minimize}\quad & \|Ax-b\|_2^2+\delta\|\Delta x\|_2^2+\eta\|x\|_2^2.
\end{aligned}
\]
这里，参数 $\delta\geq 0$ 控制了逼近解的光滑度，而参数 $\eta\geq 0$ 则用来控制其规模。

\noindent\textbf{例 6.3}\quad
最优输入设计。考虑一个动态系统，其输入标量序列为 $u(0),u(1),\cdots,u(N)$，输出标量序列为 $y(0),y(1),\cdots,y(N)$，它们通过卷积相关联：
\[
y(t)=\sum_{\tau=0}^t h(\tau)u(t-\tau),\quad t=0,1,\cdots,N.
\]
序列 $h(0),h(1),\cdots,h(N)$ 称为卷积核或系统的脉冲响应。

我们的目标是选择输入序列 $u$ 以达到一些目标：

跟踪输出：首要的目标是输出 $y$ 需要跟踪给定目标或参考信号 $y_{\mathrm{des}}$。我们用二次函数度量输出的跟踪误差
\[
J_{\mathrm{track}}=\frac{1}{N+1}\sum_{t=0}^N (y(t)-y_{\mathrm{des}}(t))^2.
\]

\begin{itemize}
\item 小的输入：输入不应该很大。我们用二次函数度量输入的幅值
\[
J_{\mathrm{mag}}=\frac{1}{N+1}\sum_{t=0}^N u(t)^2.
\]
\item 小的输入变化：输入不应当快速地变化。我们用二次函数来度量输入变化的幅值
\[
J_{\mathrm{der}}=\frac{1}{N}\sum_{t=0}^{N-1}(u(t+1)-u(t))^2.
\]
\end{itemize}

通过极小化加权和
\[
J_{\mathrm{track}}+\delta J_{\mathrm{der}}+\eta J_{\mathrm{mag}},
\]
我们可以在三个目标间进行权衡，其中 $\delta>0,\eta>0$。

这里，我们考虑一个特定的例子，$N=200$，脉冲响应为
\[
h(t)=\frac{1}{9}(0.9)^t(1-0.4\cos(2t)).
\]
图 6.6 显示了对于三组不同的正则化参数 $\delta$ 和 $\eta$ 的（关于目标轨迹 $y_{\mathrm{des}}$ 的）最优输入和相应的输出。第一行显示了对于 $\delta=0,\eta=0.005$ 的最优输入和相应的输出。在这种情况下，我们对输入的幅值进行了正则化，但没有考虑其变化。此时，跟踪良好（即 $J_{\mathrm{track}}$ 很小），而要求的输入较大，变化较快。第二行显示了对应于 $\delta=0,\eta=0.05$ 的情况。在此情况下，我们有更多的幅值正则化，但仍然没有对 $u$ 的变化进行正则化。相应的输入确实变小了，当然付出了较大的跟踪误差。最下面一行显示了 $\delta=0.3,\eta=0.05$ 时的结果。在这种情况下，我们对变化量添加了正则化。输入的变化显著地减少了，并且没有给输出带来太多的跟踪误差。

\noindent\textbf{$\ell_1$-范数正则化}\quad
$\ell_1$-范数的正则化可以用作求取稀疏解的启发式算法。例如，考虑问题
\begin{equation}
\begin{aligned}
\text{minimize}\quad & \|Ax-b\|_2+\gamma\|x\|_1,
\end{aligned}
\tag{6.11}
\end{equation}
其中残差用 Euclid 范数度量而正则化则由 $\ell_1$-范数进行。通过调整参数 $\gamma$，我们可以在 $\|Ax-b\|_2$ 和 $\|x\|_1$ 之间的最优权衡曲线上移动，由此得到 $\|Ax-b\|_2$ 和向量 $x$ 的稀疏性或基数 $\mathbf{card}(x)$（即非零元素个数）之间最优权衡曲线的近似。问题 (6.11) 可以被改写为二阶锥规划并得到求解。

\noindent\textbf{例 6.4}\quad
回归量选择问题。给定一个矩阵 $A\in\mathbb{R}^{n\times m}$，其列为潜在的回归向量，向量 $b\in\mathbb{R}^n$ 将由 $A$ 的 $k<m$ 列的线性组合进行拟合。问题是如何选择 $k$ 个回归向量并确定相应的系数。我们可以将这个问题表述为
\[
\begin{aligned}
\text{minimize}\quad & \|Ax-b\|_2\\
\text{subject to}\quad & \mathbf{card}(x)\leq k.
\end{aligned}
\]
大体上，这是一个困难的组合问题。

一个直接的解法是检查所有含有 $k$ 个非零分量的 $x$ 的可能稀疏形式。对于给定的稀疏形式，我们可以通过求解最小二乘问题得到最优的 $\tilde{x}$，即极小化 $\|\tilde{A}\tilde{x}-b\|_2$，其中 $\tilde{A}$ 是 $A$ 的子矩阵，对应于相应的稀疏形式，而 $\tilde{x}$ 是由 $x$ 非零分量组成的子向量。对于所有 $n!/(k!(n-k)!)$ 种有 $k$ 个非零量的稀疏形式，我们都需要进行求解。

一个较好的启发式方法是对于不同的 $\gamma$ 求解问题 (6.11)，寻找能够得到 $\mathbf{card}(x)=k$ 的最小的 $\gamma$。然后，我们固定这个稀疏模式，得到使 $\|Ax-b\|_2$ 最小的 $x$。

图 6.7 显示了 $A\in\mathbb{R}^{10\times 20}$，$x\in\mathbb{R}^{20}$，$b\in\mathbb{R}^{10}$ 的一个数值算例。

虚线上的圆圈是 $\mathbf{card}(x)$（纵坐标）和残差 $\|Ax-b\|_2$（横坐标）的权衡曲线上的（全局）Pareto 最优值。对于每个 $k$，Pareto 最优解通过枚举所有可能的有 $k$ 个非零分量的稀疏形式得到，如前所述。实线上的圆圈是通过启发式方法得到的，利用了不同 $\gamma$ 值对应的问题 (6.11) 的解的稀疏形式。注意到，对于 $\mathbf{card}(x)=1$，启发式算法事实上找到了全局最优解。

这个想法还将出现在基筛选问题中（$\S 6.5.4$）。

\subsection*{6.3.3\quad 重构、光滑与去除噪声}

本节已描述了双准则逼近问题的一个重要的特殊实例。我们还将给出一些例子，显示不同正则化方法的性能。在重构问题中，我们从由向量 $x\in\mathbb{R}^n$ 表示的信号出发。$x_i$ 的系数对应于某些时间函数在平均分布的点上的值（在信号处理领域称为采样）。通常假设信号不会过快地变化，也就意味着通常有 $x_i\approx x_{i+1}$。（本节考虑一维信号，例如声音信号，但是同样的想法可以应用到两维或更高维的信号中，例如，图像或影像。）信号 $x$ 被加性噪声 $v$ 所污染，即
\[
x_{\mathrm{cor}}=x+v.
\]
有很多种方法可以用来对噪声进行建模，但这里，我们简单地假设它是未知、小值的，并且快速变化（不像信号）。目标是在给定受污染信号 $x_{\mathrm{cor}}$ 的情况下，构建对原始信号 $x$ 的估计值 $\hat{x}$。这一过程称为信号重构（因为我们试图从被污染的信号中构建出原信号），或者去除噪声（因为我们试图将噪声从被污染的信号中去除）。多数重构方法最终可以视作将某些光滑运算作用在 $x_{\mathrm{cor}}$ 上以得到 $\hat{x}$，因此这一过程也称为光滑化。

重构问题的一个简单方式是双准则问题
\begin{equation}
\begin{aligned}
\text{minimize（关于 $\mathbb{R}_+^2$）}\quad & (\|\hat{x}-x_{\mathrm{cor}}\|_2,\phi(\hat{x})),
\end{aligned}
\tag{6.12}
\end{equation}
其中 $\hat{x}$ 是变量而 $x_{\mathrm{cor}}$ 为问题参数。函数 $\phi:\mathbb{R}^n\to\mathbb{R}$ 是凸的，称为正则化函数或光滑目标。这一目标被用来度量估计值 $\hat{x}$ 的粗糙度，或光滑度的缺失。重构问题 (6.12) 寻求信号，以接近（$\ell_2$-范数意义下）被污染信号并且光滑，即 $\phi(\hat{x})$ 较小。重构问题 (6.12) 是一个凸的双准则问题。我们可以通过标量化和求解（标量）凸优化问题找到 Pareto 最优解。

\noindent\textbf{二次光滑}\quad
最简单的重构方法是使用二次光滑函数
\[
\phi_{\mathrm{quad}}(x)=\sum_{i=1}^{n-1}(x_{i+1}-x_i)^2=\|Dx\|_2^2,
\]
其中，$D\in\mathbb{R}^{(n-1)\times n}$ 为双对角矩阵
\[
D=
\begin{bmatrix}
-1 & 1 & 0 & \cdots & 0 & 0 & 0\\
0 & -1 & 1 & \cdots & 0 & 0 & 0\\
\vdots & \vdots & \vdots & & \vdots & \vdots & \vdots\\
0 & 0 & 0 & \cdots & -1 & 1 & 0\\
0 & 0 & 0 & \cdots & 0 & -1 & 1
\end{bmatrix}.
\]
我们可以通过极小化
\[
\|\hat{x}-x_{\mathrm{cor}}\|_2^2+\delta\|D\hat{x}\|_2^2
\]
得到 $\|\hat{x}-x_{\mathrm{cor}}\|_2$ 和 $\|D\hat{x}\|_2$ 之间的最优权衡，其中 $\delta>0$ 参数化了最优权衡曲线。这个二次问题的解
\[
\hat{x}=(I+\delta D^T D)^{-1}x_{\mathrm{cor}}
\]
可以非常高效地求得，因为 $I+\delta D^TD$ 是三对角的；参见附录 C。

\noindent\textbf{二次光滑的例子}\quad
图 6.8 显示了信号 $x\in\mathbb{R}^{4000}$（第一行）和被污染的信号 $x_{\mathrm{cor}}$（第二行）。目标 $\|\hat{x}-x_{\mathrm{cor}}\|_2$ 和 $\|D\hat{x}\|_2$ 之间的最优权衡曲线由图 6.9 给出。在权衡曲线的左端点表示 $\hat{x}=x_{\mathrm{cor}}$ 和目标函数数值 $\|Dx_{\mathrm{cor}}\|_2=4.4$。右端点表示 $\hat{x}=0$ 及 $\|\hat{x}-x_{\mathrm{cor}}\|_2=\|x_{\mathrm{cor}}\|_2=16.2$。注意，权衡曲线在 $\|\hat{x}-x_{\mathrm{cor}}\|_2\approx 3$ 附近有一个明显的拐点。

图 6.10 显示了最优权衡曲线上的三个光滑化信号，对应于 $\|\hat{x}-x_{\mathrm{cor}}\|_2=8$（顶部），$3$（中间）和 $1$（底部）。将重构信号和原始信号 $x$ 进行对比，我们可以看出最优的重构结果在 $\|\hat{x}-x_{\mathrm{cor}}\|_2=3$ 时得到，正是曲线的拐点。$\|\hat{x}-x_{\mathrm{cor}}\|_2$ 的值更大时，结果过于光滑；而值更小时，所得结果的光滑性不足。

\noindent\textbf{总变差重构}\quad
当原始信号非常光滑而噪声变化很快时，简单的二次光滑可以作为良好的重构方法。但是显然，原始信号的任意快速变化都会被二次光滑方法所减弱或者移除。本节将描述一种重构方法，可以去除大部分的噪声，同时仍然保留原始信号偶尔的快速变化。这一方法基于光滑函数
\[
\phi_{\mathrm{tv}}(\hat{x})=\sum_{i=1}^{n-1}|\hat{x}_{i+1}-\hat{x}_i|=\|D\hat{x}\|_1,
\]
称为 $x\in\mathbb{R}^n$ 的总变差。如同二次光滑性的度量 $\phi_{\mathrm{quad}}$，总变差函数对于快速变化的 $\hat{x}$ 给予大的值。但是，总变差度量对于大的 $|x_{i+1}-x_i|$ 给予相对较小的惩罚。

\noindent\textbf{总变差重构的例子}\quad
图 6.11 显示了信号 $x\in\mathbb{R}^{2000}$（顶部）及受噪声污染的信号 $x_{\mathrm{cor}}$。信号大部分光滑，但有一些快速的变化或值上的跳跃；噪声则是快速变化的。

我们首先使用二次光滑。图 6.12 显示了处于 $\|D\hat{x}\|_2$ 和 $\|\hat{x}-x_{\mathrm{cor}}\|_2$ 之间最优权衡曲线上的三个光滑信号。在前两个信号中，原始信号的快速变化也被光滑掉了。第三个信号更好地保留了信号的阶跃边缘，但也显著地遗留了一些噪声。

现在，我们显示总变差重构的效果。图 6.13 显示了 $\|D\hat{x}\|_1$ 和 $\|\hat{x}-x_{\mathrm{cor}}\|_2$ 之间的最优权衡曲线。图 6.14 显示了处于最优权衡曲线上的重构信号，分别是 $\|D\hat{x}\|_1=5$（顶部），$\|D\hat{x}\|_1=8$（中间）和 $\|D\hat{x}\|_1=10$（底部）。我们观察到，不像二次光滑，总变差重构保留了信号的尖锐过渡。底部的信号没有给出足够的噪声衰减，而顶部的则去除了信号的某些缓慢变化的部分。需要注意，不像二次光滑，在总变差重构中，信号的尖锐变化被保留了下来。

\section*{6.4\quad 鲁棒逼近}

\subsection*{6.4.1\quad 随机鲁棒逼近}

我们研究的逼近问题的基本目标为 $\|Ax-b\|$，但是也希望考虑数据矩阵 $A$ 的一些不确定性和可能的变化。（同样的想法可以扩展到 $A$ 和 $b$ 都有不确定性的情况。）本节将考虑一些统计模型以处理 $A$ 的变化。

我们假设 $A$ 是在 $\mathbb{R}^{m\times n}$ 中取值的随机变量，其均值为 $\bar{A}$，因此，我们可以将 $A$ 描述为
\[
A=\bar{A}+U,
\]
其中 $U$ 为零均值随机矩阵。这里，常值矩阵 $\bar{A}$ 给出了 $A$ 的平均值而 $U$ 描述了其统计变化。

很自然地，可以用 $\|Ax-b\|$ 的期望值作为目标函数
\begin{equation}
\begin{aligned}
\text{minimize}\quad & \mathbf{E}\|Ax-b\|.
\end{aligned}
\tag{6.13}
\end{equation}
我们称这一问题为随机鲁棒逼近问题。这总是凸优化问题，但是通常并不容易处理，因为在大多数情况下，这个目标函数及其导数都非常难以计算。

可解的随机鲁棒逼近问题的一个简单例子是假设 $A$ 仅有有限个可能值，即
\[
\mathbf{prob}(A=A_i)=p_i,\quad i=1,\cdots,k,
\]
其中 $A_i\in\mathbb{R}^{m\times n}$，$\mathbf{1}^Tp=1$，$p\succeq 0$。在这种情况下，问题 (6.13) 具有形式
\[
\begin{aligned}
\text{minimize}\quad & p_1\|A_1x-b\|+\cdots+p_k\|A_kx-b\|,
\end{aligned}
\]
这常称为范数和问题。它可以被表述为
\[
\begin{aligned}
\text{minimize}\quad & p^Tt\\
\text{subject to}\quad & \|A_ix-b\|\leq t_i,\quad i=1,\cdots,k,
\end{aligned}
\]
其中变量为 $x\in\mathbb{R}^n$ 和 $t\in\mathbb{R}^k$。如果范数为 Euclid 范数，范数和问题是一个二阶锥规划（SOCP）。如果范数为 $\ell_1$-或 $\ell_\infty$-范数，范数和问题可以被表述为线性规划；参见习题 6.8。

易于处理随机鲁棒逼近问题 (6.13) 的一些变形。作为一个例子，考虑随机鲁棒最小二乘问题
\[
\begin{aligned}
\text{minimize}\quad & \mathbf{E}\|Ax-b\|_2^2,
\end{aligned}
\]
其中的范数为 Euclid 范数。我们可以将目标函数表示为
\[
\begin{aligned}
\mathbf{E}\|Ax-b\|_2^2
&=\mathbf{E}(\bar{A}x-b+Ux)^T(\bar{A}x-b+Ux)\\
&=(\bar{A}x-b)^T(\bar{A}x-b)+\mathbf{E}x^TU^TUx\\
&=\|\bar{A}x-b\|_2^2+x^TPx,
\end{aligned}
\]
其中 $P=\mathbf{E}U^TU$。因此，随机鲁棒逼近问题具有正则化最小二乘的形式
\[
\begin{aligned}
\text{minimize}\quad & \|\bar{A}x-b\|_2^2+\|P^{1/2}x\|_2^2,
\end{aligned}
\]
其解为
\[
x=(\bar{A}^T\bar{A}+P)^{-1}\bar{A}^Tb.
\]
这一结果非常有意义：当矩阵 $A$ 发生变化时，$x$ 越大，$Ax$ 的变化越大，Jensen 不等式告诉我们 $Ax$ 的变化会增加 $\|Ax-b\|_2$ 的平均值。因此，我们需要在使得 $\bar{A}x-b$ 较小和期望 $x$ 较小（以使得 $Ax$ 的变化较小）之间进行权衡；这实际上也是正则化的基本想法。

这一结果给出了 Tikhonov 正则化最小二乘问题 (6.10) 的另一种解释，即将其看成一种考虑了矩阵 $A$ 的可能变化的鲁棒最小二乘问题。Tikhonov 正则化最小二乘问题 (6.10) 的最优解极小化了 $\mathbf{E}\|(A+U)x-b\|^2$，其中 $U_{ij}$ 为零均值、不相关的随机变量，其方差为 $\delta/m$（并且，此处的 $A$ 是确定的）。

\subsection*{6.4.2\quad 最坏情况鲁棒逼近}

有可能采用基于集合的、最坏情况方法对矩阵 $A$ 的变化进行建模。我们用 $A$ 的可能值集合
\[
A\in\mathcal{A}\subseteq\mathbb{R}^{m\times n}
\]
描述其不确定性，这里我们假设这个集合是非空和有界的。我们定义候选的逼近解 $x\in\mathbb{R}^n$ 的最坏误差为
\[
e_{\mathrm{wc}}(x)=\sup\{\|Ax-b\|\mid A\in\mathcal{A}\},
\]
它总是 $x$ 的凸函数。（最坏情况）鲁棒逼近问题是极小化最坏情况的误差：
\begin{equation}
\begin{aligned}
\text{minimize}\quad & e_{\mathrm{wc}}(x)=\sup\{\|Ax-b\|\mid A\in\mathcal{A}\},
\end{aligned}
\tag{6.14}
\end{equation}
其中的变量为 $x$，问题数据为 $b$ 和集合 $\mathcal{A}$。当 $\mathcal{A}$ 是单点集（即 $\mathcal{A}=\{A\}$）时，鲁棒逼近问题 (6.14) 退化为基本的范数逼近问题 (6.1)。鲁棒逼近问题总是凸优化问题；但是它的求解难度与所用的范数及对不确定性集合 $\mathcal{A}$ 的描述相关。

\noindent\textbf{例 6.5}\quad
随机与最坏情况鲁棒逼近的比较。

为说明随机和最坏情况鲁棒逼近问题的差异，我们考虑最小二乘问题
\[
\begin{aligned}
\text{minimize}\quad & \|A(u)x-b\|_2^2,
\end{aligned}
\]
其中 $u\in\mathbb{R}$ 为不确定的参数，而 $A(u)=A_0+uA_1$。我们考虑这一问题的一个特定实例：$A(u)\in\mathbb{R}^{20\times 10}$，$\|A_0\|=10$，$\|A_1\|=1$ 而 $u$ 的变化区间为 $[-1,1]$。（因此，粗略地，矩阵 $A$ 的变化幅度在 $\pm 10\%$ 附近。）

我们找到了三个近似解：
\begin{itemize}
\item 名义最优。假设 $A(u)$ 的名义值为 $A_0$，我们可以找到最优解 $x_{\mathrm{nom}}$。
\item 随机鲁棒逼近。假设参数 $u$ 正态分布于 $[-1,1]$，我们可以找到 $x_{\mathrm{stoch}}$，它极小化了 $\mathbf{E}\|A(u)x-b\|_2^2$。
\item 最坏情况鲁棒逼近。我们可以找到 $x_{\mathrm{wc}}$，它极小化了
\[
\sup_{-1\leq u\leq 1}\|A(u)x-b\|_2
=\max\{\|(A_0-A_1)x-b\|_2,\|(A_0+A_1)x-b\|_2\}.
\]
\end{itemize}
对于 $x$ 的每一个结果，我们将残差视为不确定参数 $u$ 的函数并在图 6.15 中绘出 $r(u)=\|A(u)x-b\|_2$。这些图像显示了近似解对于参数 $u$ 变化有多么敏感。名义解在 $u=0$ 时达到最小的残差，但是对参数的变化相当敏感：$u$ 偏离 $0$ 趋向 $-1$ 或 $1$ 时，它给出了非常大的残差。最坏情况的解则在 $u=0$ 时有较大的残差，但当 $u$ 在区间 $[-1,1]$ 上变化时，残差并没有太大的增长。随机鲁棒逼近解介于两者之间。

鲁棒逼近问题 (6.14) 出现在很多领域和应用中。在估计范畴，集合 $\mathcal{A}$ 给出了待估计向量和测量向量之间线性关系的不确定性。有时，模型 $y=Ax+v$ 中的噪声项 $v$ 称为加性噪声或加性误差，因为它叠加在“理想”测量值 $Ax$ 上。相对应地，$A$ 的变化称为乘性误差，因为它与变量 $x$ 相乘。

在最优设计领域，变化可以表示设计变量 $x$ 和结果向量 $Ax$ 之间的线性方程的不确定性（比如，在制造过程中产生）。于是，鲁棒逼近问题 (6.14) 可以解释为鲁棒设计问题：寻找设计变量 $x$ 以极小化所有 $A$ 的可能值中 $Ax$ 和 $b$ 最坏的不匹配情况。

\noindent\textbf{有限集合}\quad
此处，我们有 $\mathcal{A}=\{A_1,\cdots,A_k\}$，其鲁棒逼近问题是
\[
\begin{aligned}
\text{minimize}\quad & \max_{i=1,\cdots,k}\|A_ix-b\|.
\end{aligned}
\]
这一问题等价于关于多面体集合 $\mathcal{A}=\mathbf{conv}\{A_1,\cdots,A_k\}$ 的鲁棒逼近问题：
\[
\begin{aligned}
\text{minimize}\quad & \sup\{\|Ax-b\|\mid A\in\mathbf{conv}\{A_1,\cdots,A_k\}\}.
\end{aligned}
\]
我们可以将这个问题描述为下述上境图形式
\[
\begin{aligned}
\text{minimize}\quad & t\\
\text{subject to}\quad & \|A_ix-b\|\leq t,\quad i=1,\cdots,k,
\end{aligned}
\]
根据使用的范数，这一问题有不同的解法。如果范数为 Euclid 范数，这是一个 SOCP。如果范数为 $\ell_1$-或 $\ell_\infty$-范数，我们可以将其表述为线性规划。

\noindent\textbf{有界范数误差}\quad
这里的不确定性集合 $\mathcal{A}$ 是一个范数球，$\mathcal{A}=\{\bar{A}+U\mid \|U\|\leq a\}$，其中 $\|\cdot\|$ 为 $\mathbb{R}^{m\times n}$ 上的一个范数。在这种情况下，我们有
\[
e_{\mathrm{wc}}(x)=\sup\{\|\bar{A}x-b+Ux\|\mid \|U\|\leq a\},
\]
这个式子需要仔细理解，因为这里出现的第一个范数定义在 $\mathbb{R}^m$ 上（用以度量误差的规模），而第二个范数定义在 $\mathbb{R}^{m\times n}$ 上（用以定义范数球 $\mathcal{A}$）。

在某些情况下，可以化简关于 $e_{\mathrm{wc}}(x)$ 的表示。例如，考虑 $\mathbb{R}^n$ 上的 Euclid 范数和相应的在 $\mathbb{R}^{m\times n}$ 上的导出范数，即最大奇异值。如果 $\bar{A}x-b\neq 0$ 且 $x\neq 0$，关于 $e_{\mathrm{wc}}(x)$ 的表达式的上确界在 $U=auv^T$ 处达到，
\[
u=\frac{\bar{A}x-b}{\|\bar{A}x-b\|_2},\qquad
v=\frac{x}{\|x\|_2},
\]
并且所得到的最坏情况误差为
\[
e_{\mathrm{wc}}(x)=\|\bar{A}x-b\|_2+a\|x\|_2.
\]
（容易验证，在 $x$ 或 $\bar{A}x-b$ 为零时，这个表达式也成立。）于是，鲁棒逼近问题 (6.14) 变为
\[
\begin{aligned}
\text{minimize}\quad & \|\bar{A}x-b\|_2+a\|x\|_2,
\end{aligned}
\]
这是一个正则化范数问题，可以作为 SOCP 得到求解：
\[
\begin{aligned}
\text{minimize}\quad & t_1+at_2\\
\text{subject to}\quad & \|\bar{A}x-b\|_2\leq t_1,\quad \|x\|_2\leq t_2.
\end{aligned}
\]
因为这个问题的解与对于某个正则化参数 $\delta$ 的正则化最小二乘问题
\[
\begin{aligned}
\text{minimize}\quad & \|\bar{A}x-b\|_2^2+\delta\|x\|_2^2
\end{aligned}
\]
的解相同，所以，我们对正则化最小二乘问题有了另一个解释，可以将其看成一个最坏情况鲁棒逼近问题。

\noindent\textbf{不确定性椭球}\quad
我们也可以通过给定每一行的可能值的椭圆来描述 $A$ 的变化：
\[
\mathcal{A}=\{[a_1\ \cdots\ a_m]^T\mid a_i\in\mathcal{E}_i,\ i=1,\cdots,m\},
\]
其中
\[
\mathcal{E}_i=\{\bar{a}_i+P_iu\mid \|u\|_2\leq 1\}.
\]
矩阵 $P_i\in\mathbb{R}^{n\times n}$ 描述了 $a_i$ 的变化。我们允许 $P_i$ 有一个非平凡的零空间，用以对 $a_i$ 的变化被约束在一个子空间的情况进行建模。作为一个极端的例子，如果 $a_i$ 没有不确定性，我们取 $P_i=0$。

应用椭圆不确定性的描述，我们可以给出每一个残差在最坏情况下的幅值的显式表达式：
\[
\begin{aligned}
\sup_{a_i\in\mathcal{E}_i}|a_i^Tx-b_i|
&=\sup\{|\bar{a}_i^Tx-b_i+(P_iu)^Tx|\mid \|u\|_2\leq 1\}\\
&=|\bar{a}_i^Tx-b_i|+\|P_i^Tx\|_2.
\end{aligned}
\]
利用这个结果，我们可以求解一些鲁棒逼近问题。例如，鲁棒 $\ell_2$-范数逼近问题
\[
\begin{aligned}
\text{minimize}\quad & e_{\mathrm{wc}}(x)=\sup\{\|Ax-b\|_2\mid a_i\in\mathcal{E}_i,\ i=1,\cdots,m\}
\end{aligned}
\]
可以简化为如下所示的 SOCP。最坏情况误差的显式表达式由
\[
e_{\mathrm{wc}}(x)=
\left(\sum_{i=1}^m\left(\sup_{a_i\in\mathcal{E}_i}|a_i^Tx-b_i|\right)^2\right)^{1/2}
=
\left(\sum_{i=1}^m\left(|\bar{a}_i^Tx-b_i|+\|P_i^Tx\|_2\right)^2\right)^{1/2}
\]
给出。为极小化 $e_{\mathrm{wc}}(x)$，我们可以求解
\[
\begin{aligned}
\text{minimize}\quad & \|t\|_2\\
\text{subject to}\quad & |\bar{a}_i^Tx-b_i|+\|P_i^Tx\|_2\leq t_i,\quad i=1,\cdots,m;
\end{aligned}
\]
这里我们引入了新的变量 $t_1,\cdots,t_m$。这一问题可以表述为
\[
\begin{aligned}
\text{minimize}\quad & \|t\|_2\\
\text{subject to}\quad & \bar{a}_i^Tx-b_i+\|P_i^Tx\|_2\leq t_i,\quad i=1,\cdots,m\\
& -\bar{a}_i^Tx+b_i+\|P_i^Tx\|_2\leq t_i,\quad i=1,\cdots,m,
\end{aligned}
\]
当转变为上境图形式时，可以转换得到 SOCP。

\noindent\textbf{线性结构的范数有界误差}\quad
作为范数有界描述 $\mathcal{A}=\{\bar{A}+U\mid \|U\|\leq a\}$ 的一个扩展，我们可以将 $\mathcal{A}$ 定义为范数球在仿射变换下的像：
\[
\mathcal{A}=\{\bar{A}+u_1A_1+u_2A_2+\cdots+u_pA_p\mid \|u\|\leq 1\},
\]
其中 $\|\cdot\|$ 为 $\mathbb{R}^p$ 上的范数，而 $p+1$ 个矩阵 $\bar{A},A_1,\cdots,A_p\in\mathbb{R}^{m\times n}$ 给定。最坏情况误差可以表示为
\[
\begin{aligned}
e_{\mathrm{wc}}(x)
&=\sup_{\|u\|\leq 1}\|(\bar{A}+u_1A_1+\cdots+u_pA_p)x-b\|\\
&=\sup_{\|u\|\leq 1}\|P(x)u+q(x)\|,
\end{aligned}
\]
其中，$P$ 和 $q$ 定义为
\[
P(x)=
\begin{bmatrix}
A_1x & A_2x & \cdots & A_px
\end{bmatrix}\in\mathbb{R}^{m\times p},
\qquad q(x)=\bar{A}x-b\in\mathbb{R}^m.
\]

作为第一个例子，我们考虑鲁棒 Chebyshev 逼近问题
\[
\begin{aligned}
\text{minimize}\quad & e_{\mathrm{wc}}(x)=\sup_{\|u\|_\infty\leq 1}\|(\bar{A}+u_1A_1+\cdots+u_pA_p)x-b\|_\infty.
\end{aligned}
\]
在这种情况下，我们可以导出最坏情况误差的显式表达式。用 $p_i(x)^T$ 表示 $P(x)$ 的第 $i$ 行，我们有
\[
\begin{aligned}
e_{\mathrm{wc}}(x)
&=\sup_{\|u\|_\infty\leq 1}\|P(x)u+q(x)\|_\infty\\
&=\max_{i=1,\cdots,m}\sup_{\|u\|_\infty\leq 1}|p_i(x)^Tu+q_i(x)|\\
&=\max_{i=1,\cdots,m}(\|p_i(x)\|_1+|q_i(x)|).
\end{aligned}
\]
因此，鲁棒 Chebyshev 逼近问题可以被表述为线性规划
\[
\begin{aligned}
\text{minimize}\quad & t\\
\text{subject to}\quad & -y_0\preceq \bar{A}x-b\preceq y_0\\
& -y_k\preceq A_kx\preceq y_k,\quad k=1,\cdots,p\\
& y_0+\sum_{k=1}^p y_k\preceq t\mathbf{1},
\end{aligned}
\]
其变量为 $x\in\mathbb{R}^n$，$y_k\in\mathbb{R}^m$，$t\in\mathbb{R}$。

作为另一个例子，我们考虑鲁棒最小二乘问题
\[
\begin{aligned}
\text{minimize}\quad & e_{\mathrm{wc}}(x)=\sup_{\|u\|_2\leq 1}\|(\bar{A}+u_1A_1+\cdots+u_pA_p)x-b\|_2.
\end{aligned}
\]
这里，我们用 Lagrange 对偶来计算 $e_{\mathrm{wc}}$。最坏误差 $e_{\mathrm{wc}}(x)$ 是下面（非凸）二次优化问题最优解的平方根
\[
\begin{aligned}
\text{maximize}\quad & \|P(x)u+q(x)\|_2^2\\
\text{subject to}\quad & u^Tu\leq 1,
\end{aligned}
\]
其中 $u$ 为变量。这一问题的 Lagrange 对偶问题可以被表述为半定规划（SDP）
\begin{equation}
\begin{aligned}
\text{minimize}\quad & t+\lambda\\
\text{subject to}\quad &
\begin{bmatrix}
I & P(x) & q(x)\\
P(x)^T & \lambda I & 0\\
q(x)^T & 0 & t
\end{bmatrix}\succeq 0,
\end{aligned}
\tag{6.15}
\end{equation}
其变量为 $t,\lambda\in\mathbb{R}$。并且，如 $\S 5.2$ 和 $\S B.1$ 所述（并在 $\S B.4$ 中证明），对于这组原-对偶问题，强对偶成立。换言之，对于固定的 $x$，我们可以通过求解关于变量 $t$ 和 $\lambda$ 的 SDP (6.15) 来计算 $e_{\mathrm{wc}}(x)^2$。联合优化 $t,\lambda$ 及 $x$ 等同于极小化 $e_{\mathrm{wc}}(x)^2$。我们可以断言鲁棒最小二乘问题等价于以 $x,\lambda,t$ 为变量的 SDP (6.15)。

\noindent\textbf{例 6.6}\quad
最坏情况鲁棒、Tikhonov 正则化和名义最小二乘解的比较。我们考虑一个鲁棒逼近问题的实例
\begin{equation}
\begin{aligned}
\text{minimize}\quad & \sup_{\|u\|_2\leq 1}\|(\bar{A}+u_1A_1+u_2A_2)x-b\|_2,
\end{aligned}
\tag{6.16}
\end{equation}
其维数为 $m=50,n=20$。矩阵 $\bar{A}$ 范数为 $10$，矩阵 $A_1$ 和 $A_2$ 范数为 $1$，所以，粗略地讲，矩阵 $A$ 的变化幅度在 $10\%$ 左右。参数 $u_1$ 和 $u_2$ 的不确定性分布在 $\mathbb{R}^2$ 的单位圆盘上。

我们计算鲁棒最小二乘问题 (6.16) 的最优解 $x_{\mathrm{rls}}$，也计算名义最小二乘问题的解 $x_{\mathrm{ls}}$（即假设 $u=0$），以及 $\delta=1$ 时的 Tikhonov 正则化的解 $x_{\mathrm{tik}}$。

为显示各个近似解对于参数 $u$ 的敏感性，我们产生均匀分布于单位圆盘上的 $10^5$ 个参数向量，并对每一个参数值计算残差
\[
\|(A_0+u_1A_1+u_2A_2)x-b\|_2.
\]
残差的分别显示于图 6.16。

我们可以观察到一些结果。首先，名义最小二乘解的残差散布广泛，从 $0.52$ 附近的最小值直到 $4.9$ 附近的最大值。特别地，最小二乘解对于参数的变化非常敏感。相对应地，鲁棒最小二乘和 Tikhonov 正则化的解展现了对于在整个单位圆上参数的不确定性的相当小的变化。例如，当参数处于单位圆盘时，鲁棒最小二乘解的残差都在 $2.0$ 和 $2.6$ 之间。

\section*{6.5\quad 函数拟合与插值}

在函数拟合问题中，我们在函数的有限维子空间中，选择与给定数据或需求最符合的一个。为简单起见，我们考虑实值函数；也很容易将其思想扩展到向量函数的处理中。

\subsection*{6.5.1\quad 函数族}

我们考虑一族具有共同定义域 $\mathbf{dom}\ f_i=D$ 的函数 $f_1,\cdots,f_n:\mathbb{R}^k\to\mathbb{R}$。对于每一个 $x\in\mathbb{R}^n$，我们通过给定
\begin{equation}
f(u)=x_1f_1(u)+\cdots+x_nf_n(u)
\tag{6.17}
\end{equation}
将 $x$ 与 $f:\mathbb{R}^k\to\mathbb{R}$ 联系起来，其中 $\mathbf{dom}\ f=D$。函数族 $\{f_1,\cdots,f_n\}$ 有时也称为基函数（对于拟合问题），不论这些函数是否独立。向量 $x\in\mathbb{R}^n$ 参数化了函数子空间，是我们的优化变量，有时也称为系数向量。基函数生成了 $D$ 上的一个函数子空间 $\mathcal{F}$。

在很多应用中，基函数是根据先验知识或经验而特别选择的，用以合理地用有限维函数子空间对感兴趣的函数进行建模。在其他情况下，更多的是使用通用函数族。我们将在下面描述其中一些。

\noindent\textbf{多项式}\quad
$\mathbb{R}$ 上一个常见的函数子空间由阶次小于 $n$ 的多项式组成。最简单的基由幂组成，即 $f_i(t)=t^{i-1}$，$i=1,\cdots,n$。在很多应用中，同样的子空间可以由不同的基进行描述，例如，与某些正函数（或测量） $\phi:\mathbb{R}^n\to\mathbb{R}_+$ 正交的，阶次小于 $n$ 的多项式集合 $f_1,\cdots,f_n$，即
\[
\int f_i(t)f_j(t)\phi(t)\,dt=
\begin{cases}
1 & i=j\\
0 & i\neq j.
\end{cases}
\]
多项式的另一组常用的基是 Lagrange 基 $f_1,\cdots,f_n$，它们与离散的点 $t_1,\cdots,t_n$ 相关联，并且满足
\[
f_i(t_j)=
\begin{cases}
1 & i=j\\
0 & i\neq j.
\end{cases}
\]
我们也可以考虑 $\mathbb{R}^k$ 上的多项式，它们具有最大的总阶次限制或每个变量都有最大阶次限制。

作为一个相关的例子，我们有阶次小于 $n$ 的三角多项式，其基函数为
\[
\sin kt,\quad k=1,\cdots,n-1,\qquad \cos kt,\quad k=0,\cdots,n-1.
\]

\noindent\textbf{分片线性函数}\quad
我们从定义域 $D$ 的三角化开始讨论，其意义如下。我们有网格点集合 $g_1,\cdots,g_n\in\mathbb{R}^k$ 将 $D$ 划分为单纯形集合：
\[
D=S_1\cup\cdots\cup S_m,\qquad \mathbf{int}(S_i\cap S_j)=\emptyset \quad \text{对于 } i\neq j.
\]
每个单纯形是 $k+1$ 个网格点的凸包，我们还要求网格点是它所在的单纯形的顶点。

给定一组三角化，我们可以构造一个分片线性（或更准确地，分片仿射）函数 $f$，其方法是将网格点赋值为 $f(g_i)=x_i$，然后将函数仿射地扩展到每个单纯形上。函数 $f$ 可以如 (6.17) 一样表示，其基函数 $f_i$ 在每个单纯形上是仿射的并且由条件
\[
f_i(g_j)=
\begin{cases}
1 & i=j\\
0 & i\neq j
\end{cases}
\]
所定义。通过这种构造所得的函数是连续的。

图 6.17 显示了 $k=2$ 的一个例子。

\noindent\textbf{分片多项式和样条}\quad
三角化定义域上分片仿射函数的想法可以很容易地扩展到分片的多项式和其他函数上。

分片多项式是定义在三角化的单纯形上的（具有某些最大阶次的）多项式，它们是连续的，即多项式在单纯形的边界上相等。通过进一步的限制，可以要求分片多项式具有一定程度的连续导数；我们可以定义多种样条函数类。图 6.18 中的例子显示了立方样条，即 $\mathbb{R}$ 上的 3 阶分片多项式，它具有连续的一阶和二阶导数。

\subsection*{6.5.2\quad 约束}

本节增加对函数 $f$ 的某些约束，因此也增加了对变量 $x\in\mathbb{R}^n$ 的约束。

\noindent\textbf{函数插值与不等式}\quad
令 $v$ 为 $D$ 中的一个点。$f$ 在 $v$ 处的值
\[
f(v)=\sum_{i=1}^n x_if_i(v)
\]
是 $x$ 的线性函数。因此，插值条件
\[
f(v_j)=z_j,\quad j=1,\cdots,m
\]
要求函数在特定点 $v_j\in D$ 上具有值 $z_j\in\mathbb{R}$，这些条件构成了 $x$ 的一组线性方程。更一般地，关于给定点的函数值的不等式，如 $l\leq f(v)\leq u$，是变量 $x$ 的线性不等式。还有其他很多关于 $f$ 的有意义的凸约束（因此，也是对 $x$ 的凸约束），它们考虑了在有限点集 $v_1,\cdots,v_N$ 上的函数值。例如，Lipschitz 约束
\[
|f(v_j)-f(v_k)|\leq L\|v_j-v_k\|,\quad j,k=1,\cdots,m
\]
构成了 $x$ 的线性不等式组。

我们也可以对无穷个点的函数值加以不等式约束。作为一个例子，考虑非负约束
\[
f(u)\geq 0,\quad \forall u\in D.
\]
这是 $x$ 上的一个凸约束（因为它是无穷个半空间的交集），但是有可能会导致一个难解的问题，除非利用函数的特殊结构。一个简单的例子是函数为分片线性的时候。在这种情况下，如果函数值在网格点上非负，那么函数在各处都是非负的，因此我们可以得到一个简单的（有限的）线性不等式组。

作为一个更有意义的例子，考虑函数为 $\mathbb{R}$ 上最大阶次为偶数 $2k$（即 $n=2k+1$）的多项式，并且 $D=\mathbb{R}$ 的情况。如第 59 页的习题 2.37 所示，非负约束
\[
p(u)=x_1+x_2u+\cdots+x_{2k+1}u^{2k}\geq 0,\quad \forall u\in\mathbb{R}
\]
等价于
\[
x_i=\sum_{m+n=i+1}Y_{mn},\quad i=1,\cdots,2k+1,\qquad Y\succeq 0,
\]
其中 $Y\in\mathbb{S}^{k+1}$ 为辅助变量。

\noindent\textbf{导数约束}\quad
设基函数 $f_i$ 在 $v\in D$ 处可微。梯度
\[
\nabla f(v)=\sum_{i=1}^n x_i\nabla f_i(v)
\]
是 $x$ 的线性函数，因此，$v$ 处 $f$ 的导数的插值条件将简化为 $x$ 的线性等式约束。要求 $v$ 处梯度的范数不超过给定的限制，
\[
\|\nabla f(v)\|=\left\|\sum_{i=1}^n x_i\nabla f_i(v)\right\|\leq M,
\]
这是 $x$ 上的一个凸约束。同样的想法可以扩展到更高的导数。例如，如果 $f$ 在 $v$ 处是二阶可微的，那么要求
\[
lI\preceq \nabla^2 f(v)\preceq uI,
\]
这是 $x$ 的线性矩阵不等式，因此是凸的。

我们也可以在无限个点上对导数施加约束。例如，我们可以要求 $f$ 是单调的：
\[
f(u)\geq f(v) \quad \text{对于所有 } u,v\in D,\ u\succeq v.
\]
这是 $x$ 的凸约束，但是除非特殊的情况，所得到的问题并不容易求解。例如，当 $f$ 是分片仿射时，单调性约束等同于在每个单纯形内的条件 $\nabla f(v)\succeq 0$。因为导数是网格点上函数值的线性函数，这个条件会得到简单的（有限的）线性不等式组。

作为另一个例子，我们可以要求函数是凸的，即满足
\[
f((u+v)/2)\leq (f(u)+f(v))/2,\quad \forall u,v\in D
\]
（当 $f$ 连续时，这个条件足以保证凸性）。这是凸约束，在某些情况下，它有易于处理的表示。一个显然的例子是当 $f$ 为二次型时，在此类情况下，凸约束退化为对 $f$ 的二次部分的非负约束，这是线性矩阵不等式（LMI）。另一个由凸性约束得到易解问题的例子将在 $\S 6.5.5$ 中进行详细描述。

\noindent\textbf{积分约束}\quad
函数子空间中的任意线性泛函 $\mathcal{L}$ 都可以表示为 $x$ 的线性函数，即，我们有 $\mathcal{L}(f)=c^Tx$。对于 $f$（或其导数）在一个点的度量只是一个简单的情况。例如，线性泛函
\[
\mathcal{L}(f)=\int_D \phi(u)f(u)\,du,
\]
这里的 $\phi:\mathbb{R}^k\to\mathbb{R}$ 可以表示为 $\mathcal{L}(f)=c^Tx$，其中
\[
c_i=\int_D \phi(u)f_i(u)\,du.
\]
所以，形如 $\mathcal{L}(f)=a$ 的约束是对 $x$ 的线性等式约束。此类约束的一个例子是矩约束
\[
\int_D t^m f(t)\,dt=a
\]
（其中 $f:\mathbb{R}\to\mathbb{R}$）。

\subsection*{6.5.3\quad 拟合与插值问题}

\noindent\textbf{函数的最小范数拟合}\quad
在拟合问题中，我们寻找函数 $f\in\mathcal{F}$ 尽量准确地匹配给定数据
\[
(u_1,y_1),\quad \cdots,\quad (u_m,y_m),
\]
其中 $u_i\in D$，$y_i\in\mathbb{R}$。例如，在最小二乘拟合中，我们考虑如下问题
\[
\begin{aligned}
\text{minimize}\quad & \sum_{i=1}^m(f(u_i)-y_i)^2,
\end{aligned}
\]
这是关于变量 $x$ 的简单的最小二乘问题。我们可以对其添加各种约束，例如 $f$ 在一些点上必须满足的线性不等式、关于 $f$ 的导数的约束、单调性约束或矩约束。

\noindent\textbf{例 6.7}\quad
多项式拟合。对于给定的数据 $u_1,\cdots,u_m\in\mathbb{R}$ 和 $v_1,\cdots,v_m\in\mathbb{R}$，我们希望用具有如下形式的多项式近似地拟合数据，
\[
p(u)=x_1+x_2u+\cdots+x_nu^{n-1}.
\]
对于每个 $x$，我们构造误差向量
\[
e=(p(u_1)-v_1,\cdots,p(u_m)-v_m).
\]
为寻找极小化误差范数的多项式，我们求解关于变量 $x\in\mathbb{R}^n$ 的范数逼近问题
\[
\begin{aligned}
\text{minimize}\quad & \|e\|=\|Ax-v\|,
\end{aligned}
\]
其中 $A_{ij}=u_i^{j-1}$，$i=1,\cdots,m$，$j=1,\cdots,n$。

图 6.19 显示了对 $m=40$ 个数据点且 $n=6$（即多项式最高阶次为 5）的例子进行 $\ell_2$-和 $\ell_\infty$-范数拟合的情况。

\noindent\textbf{例 6.8}\quad
样条拟合。图 6.20 显示了对于例 6.7 中数据，三次样条的两个最佳拟合效果。区间 $[-1,1]$ 被等分为三个区间，我们考虑最大阶次为 3 的分片多项式，同时保证一阶和二阶导数的连续性。这个函数子空间的维数为 6，与例 6.7 中考虑的、最大阶次为 5 的多项式空间维数相同。

在最简单的函数拟合的形式中，我们有 $m\gg n$，即数据点个数远大于函数子空间的维数。光滑性是自动完成的，因为子空间中的所有成员都是光滑的。

\noindent\textbf{最小范数插值}\quad
在函数拟合的另一个变形中，我们有少于函数子空间维数的数据点。在最简单的情况中，我们所选的函数必须满足插值条件
\[
f(u_i)=y_i,\quad i=1,\cdots,m,
\]
这是关于 $x$ 的线性等式约束。在所有满足插值条件的函数中，我们或许会选择最为光滑或最小的一个。这导致了最小范数问题。

在最广义的函数逼近问题中，我们可以在一些代表潜在函数的先验知识的凸约束下，优化目标（比如对于误差 $e$ 的某些度量）。

\noindent\textbf{插值，外推及区间限定}\quad
通过对原始数据集外一点 $v$ 的最优函数拟合 $\hat{f}$ 进行评价，我们可以猜想潜在函数在 $v$ 处的值是多少。当 $v$ 处于给定数据点之间或附近（即 $v\in\mathbf{conv}\{v_1,\cdots,v_m\}$）时，这称为内插，否则称为外推。

我们也可以通过在一些约束下，极大和极小化（线性函数） $f(v)$，从而得到函数值 $f(v)$ 的可能区间。我们可以利用函数拟合来帮助辨别错误的数据或野值。例如，我们可以利用 $\ell_1$-范数逼近来寻找具有大的误差的数据点。

\subsection*{6.5.4\quad 稀疏描述与基筛选}

基筛选中所考虑的基函数的数量非常大，目的是寻找较少数目的基函数，使其线性组合很好地拟合给定数据。（在此类语境下，函数族线性相关，有时称为过完全基或字典。）这一过程称为基筛选，因为我们将从给定的过完全基中，挑选（远小于给定的数量的）基函数用以对数据进行建模。

因此，我们寻求函数 $f\in\mathcal{F}$ 使得能够由一个稀疏的系数向量 $x$，即 $\mathbf{card}(x)$ 很小，来很好地拟合数据
\[
f(u_i)\approx y_i,\quad i=1,\cdots,m.
\]
在此情况下，我们称
\[
f=x_1f_1+\cdots+x_nf_n=\sum_{i\in\mathcal{B}}x_if_i
\]
为数据的一个稀疏描述，其中，集合 $\mathcal{B}=\{i\mid x_i\neq 0\}$ 的元素是选中的基函数的下标集合。在数学上，基筛选等同于回归向量选择问题（参见 $\S 6.4$），但是相应的优化问题的解释（和规模）是不同的。

稀疏描述和基筛选有很多用处。它们可以用来去除噪声或光滑化，或者在信号的有效存储和传输中用来进行数据压缩。在数据压缩中，发送者和接收者均知道字典或基元素。为将信号发送给接收者，发送者首先寻找信号的稀疏表示，然后仅将（一定精度下）的非零系数发送给接收者。使用这些系数，接收者可以重构出原始信号（的一个近似）。

常用的基筛选的方法与在 $\S 6.4$ 中描述的回归向量选择问题相同，即将基于 $\ell_1$-范数的正则化视为寻找稀疏表示的一个启发式算法。我们首先求解凸问题
\begin{equation}
\begin{aligned}
\text{minimize}\quad & \sum_{i=1}^m(f(u_i)-y_i)^2+\gamma\|x\|_1,
\end{aligned}
\tag{6.18}
\end{equation}
其中 $\gamma>0$ 是一个用以权衡对数据的拟合质量和系数向量稀疏度的参数。这一问题的解可以直接拿来使用，或者再进行一步改进，用 (6.18) 的解的稀疏模式来寻找最优拟合。换言之，我们首先求解 (6.18) 来得到 $\hat{x}$。然后，令 $\mathcal{B}=\{i\mid \hat{x}_i\neq 0\}$，即非零系数对应的指标集合，并求解最小二乘问题
\[
\begin{aligned}
\text{minimize}\quad & \sum_{i=1}^m(f(u_i)-y_i)^2,
\end{aligned}
\]
其变量为 $x_i$，$i\in\mathcal{B}$，而对于 $i\notin\mathcal{B}$，有 $x_i=0$。

在基筛选和稀疏描述的应用中，经常会有一个非常大的字典，$n$ 会在 $10^4$ 的数量级上，甚至更多。为高效起见，求解 (6.18) 的算法必须利用问题的结构，这可从字典信号的结构中推知。

\noindent\textbf{通过基筛选的时间-频率分析}\quad
本节将用一个简单的例子来显示基筛选和稀疏表示。考虑 $\mathbb{R}$ 上的函数（或信号），感兴趣的范围为 $[0,1]$。我们将独立的变量视为时间，因此用 $t$（代替 $u$）来表示它。

我们首先描述字典中的基函数。每个基函数为 Gauss 正弦脉冲或 Gabor 函数，其形式为
\[
e^{-(t-\tau)^2/\sigma^2}\cos(\omega t+\phi),
\]
其中 $\sigma>0$ 给出了脉冲的宽度，$\tau$ 为脉冲的时间（的中心），$\omega\geq 0$ 为频率，而 $\phi$ 为相角。所有基函数具有同样的宽度 $\sigma=0.05$。脉冲时间和频率为
\[
\tau=0.002k,\quad k=0,\cdots,500,\qquad \omega=5k,\quad k=0,\cdots,30.
\]
对于每个时间 $\tau$，对应零频率（且相位 $\phi=0$）有 1 个基元素；对应其他 30 个频率，各有 2 个基元素（正弦和余弦，即相位 $\phi=0$ 和 $\phi=\pi/2$）。所以，总共有 $501\times 61=30561$ 个基元素。这些基元素自然地按照时间、频率和相位（正弦或余弦）进行排序，我们将其记为
\[
f_{\tau,\omega,c},\quad \tau=0,0.002,\cdots,1,\quad \omega=0,5,\cdots,150,
\]
\[
f_{\tau,\omega,s},\quad \tau=0,0.002,\cdots,1,\quad \omega=5,\cdots,150.
\]
图 6.21 显示了这些基函数中的三个（均对应时间 $\tau=0.5$）。

这个字典中的基筛选可以视为数据的时间-频率分析。如果基元素 $f_{\tau,\omega,c}$ 或 $f_{\tau,\omega,s}$ 出现在信号的稀疏表示中（即具有非零系数），那么，我们可以将其解释为数据在时间 $\tau$ 包含频率 $\omega$。

我们将利用基筛选来找到信号
\[
y(t)=a(t)\sin\hat{\theta}(t)
\]
的一个稀疏逼近，其中
\[
a(t)=1+0.5\sin(11t),\qquad \hat{\theta}(t)=30\sin(5t).
\]
（选择这个信号仅仅是因为它易于描述和显示频谱关于时间的显著变化。）我们可以将 $a(t)$ 解释为信号幅值，而将 $\hat{\theta}(t)$ 解释为总相位。我们也可以将
\[
\omega(t)=\left|\frac{d\theta}{dt}\right|=150|\cos(5t)|
\]
理解为信号在时刻 $t$ 的瞬时频率。给定数据为均匀分布于区间 $[0,1]$ 上的 501 个样本，即我们有 501 组 $(t_k,y_k)$，
\[
t_k=0.005k,\qquad y_k=y(t_k),\qquad k=0,\cdots,500.
\]
我们首先求解 $\ell_1$-范数正则化最小二乘问题 (6.18)，这里 $\gamma=1$。所得到的最优向量非常稀疏，30561 个系数中仅有 42 个非零。然后，我们用这 42 个基向量得到原始信号的最小二乘拟合。其结果 $\hat{y}$ 与原始信号 $y$ 在图 6.22 中进行了比较。顶部的图显示了逼近信号（虚线），它与原始信号 $y(t)$（实线）几乎不可分辨。底部的图显示了误差 $y(t)-\hat{y}(t)$。从图中可以清晰地看出，我们得到了具有非常好的拟合特性的逼近结果 $\hat{y}$，其相对误差为

\[
\frac{(1/501)\sum_{i=1}^{501}(y(t_i)-\hat{y}(t_i))^2}{(1/501)\sum_{i=1}^{501}y(t_i)^2}=2.6\cdot 10^{-4}.
\]
通过绘制与非零系数相对应的时间和频率模式，我们得到了原始数据的时间-频率分析。这些点及瞬时频率显示在图 6.23 中。图像表明，非零分量紧密地跟踪了瞬时频率。

\subsection*{6.5.5\quad 利用凸函数的插值}

在一些特殊的情况下，我们可以利用有限维凸优化来解决无穷维函数集合的插值问题。我们将在这一节描述一个例子。

我们从下面的问题说起：什么时候存在凸函数 $f:\mathbb{R}^k\to\mathbb{R}$，$\mathbf{dom}\ f=\mathbb{R}^k$，在给定点 $u_i\in\mathbb{R}^k$ 满足插值条件
\[
f(u_i)=y_i,\quad i=1,\cdots,m,
\]
（这里，我们不将 $f$ 限制于任何有限维函数子空间。）这个问题的答案是：当且仅当存在 $g_1,\cdots,g_m$ 使得
\begin{equation}
y_j\geq y_i+g_i^T(u_j-u_i),\quad i,j=1,\cdots,m.
\tag{6.19}
\end{equation}
为说明这点，首先设 $f$ 是凸的，$\mathbf{dom}\ f=\mathbb{R}^k$，并且 $f(u_i)=y_i$，$i=1,\cdots,m$。在每个 $u_i$，我们可以找到一个向量 $g_i$ 使得对于所有 $z$ 都有
\begin{equation}
f(z)\geq f(u_i)+g_i^T(z-u_i).
\tag{6.20}
\end{equation}
如果 $f$ 可微，我们可以取 $g_i=\nabla f(u_i)$；对于更一般的情况，我们可以通过寻找 $\mathbf{epi}\ f$ 在 $(u_i,y_i)$ 处的一个支撑超平面来构造 $g_i$。（向量 $g_i$ 称为次梯度。）将 $z=u_j$ 代入 (6.20)，我们得到 (6.19)。

反之，设 $g_1,\cdots,g_m$ 满足 (6.19)。对所有 $z$，定义 $f$ 为
\[
f(z)=\max_{i=1,\cdots,m}\left(y_i+g_i^T(z-u_i)\right).
\]
显然，$f$ 是一个（分片线性的）凸函数。不等式 (6.19) 表明，对于 $i=1,\cdots,m$ 均有 $f(u_i)=y_i$。

我们可以利用这个结果求解一些涉及（利用凸函数的）插值、逼近和定界的问题。

\noindent\textbf{对给定数据的凸函数拟合}\quad
最简单的应用可能是计算凸函数对于给定数据 $(u_i,y_i)$，$i=1,\cdots,m$ 的最小二乘拟合：
\[
\begin{aligned}
\text{minimize}\quad & \sum_{i=1}^m(y_i-f(u_i))^2\\
\text{subject to}\quad & f:\mathbb{R}^k\to\mathbb{R}\text{ 为凸函数},\quad \mathbf{dom}\ f=\mathbb{R}^k.
\end{aligned}
\]
这是一个无穷维问题，因为变量 $f$ 属于 $\mathbb{R}^k$ 上的连续、实值函数空间。利用前述结论，我们可以将问题构建为
\[
\begin{aligned}
\text{minimize}\quad & \sum_{i=1}^m(y_i-\hat{y}_i)^2\\
\text{subject to}\quad & \hat{y}_j\geq \hat{y}_i+g_i^T(u_j-u_i),\quad i,j=1,\cdots,m,
\end{aligned}
\]
这是一个关于 $\hat{y}\in\mathbb{R}^m$ 和 $g_1,\cdots,g_m\in\mathbb{R}^k$ 的二次规划。这个问题最优值为零的充要条件是，给定的数据可以被一个凸函数所插值，即存在满足 $f(u_i)=y_i$ 的凸函数。图 6.24 显示了一个例子。

\noindent\textbf{对插值凸函数的值的界定}\quad
作为另外一个简单的例子，假设给定数据 $(u_i,y_i)$，$i=1,\cdots,m$ 可以被一个凸函数所插值。我们希望确定 $f(u_0)$ 可能的取值范围，其中 $u_0$ 为 $\mathbb{R}^k$ 中的另一个点，$f$ 为给
定数据的任意一个插值函数。为得到 $f(u_0)$ 的最小可能值，我们求解线性规划（LP）
\[
\begin{aligned}
\text{minimize}\quad & y_0\\
\text{subject to}\quad & y_j\geq y_i+g_i^T(u_j-u_i),\quad i,j=0,\cdots,m,
\end{aligned}
\]
这是关于变量 $y_0\in\mathbb{R}$，$g_0,\cdots,g_m\in\mathbb{R}^k$ 的 LP。通过极大化 $y_0$（也是一个 LP），我们可以找到对于给定数据的插值凸函数的 $f(u_0)$ 的最大可能值。

\noindent\textbf{利用单调凸函数的插值}\quad
作为凸插值的一个扩展，我们可以考虑利用凸的单调非减函数进行插值。存在凸函数 $f:\mathbb{R}^k\to\mathbb{R}$，$\mathbf{dom}\ f=\mathbb{R}^k$ 满足插值条件
\[
f(u_i)=y_i,\quad i=1,\cdots,m,
\]
并且单调非减（即只要 $u\succeq v$，就有 $f(u)\geq f(v)$）的充要条件是，存在 $g_1,\cdots,g_m\in\mathbb{R}^k$ 使得
\begin{equation}
g_i\succeq 0,\quad i=1,\cdots,m,\qquad y_j\geq y_i+g_i^T(u_j-u_i),\quad i,j=1,\cdots,m.
\tag{6.21}
\end{equation}
换言之，我们向凸插值条件 (6.19) 中添加了要求次梯度非负的条件。（参见习题 6.12。）

\noindent\textbf{界定消费者的偏好}\quad
作为一个应用，我们考虑对消费者偏好的预测问题。考虑由 $n$ 种不同商品组成的货篮，货篮可以由向量 $x\in[0,1]^n$ 描述，其中 $x_i$ 表示商品 $i$ 的量。我们假设商品量是归一化了的，因此 $0\leq x_i\leq 1$，也就是说 $x_i=0$ 是货物 $i$ 的最小可能量而 $x_i=1$ 为最大可能。给定两个货篮 $x$ 和 $\tilde{x}$，消费者可以偏向于 $x$，或偏向于 $\tilde{x}$，也可以认为 $x$ 和 $\tilde{x}$ 具有等同的吸引力。我们考虑模型消费者，其选择可重复。

我们按以下方式对消费者的偏好进行建模。假设存在一个潜在的效用函数 $u:\mathbb{R}^n\to\mathbb{R}$，其定义域为 $[0,1]^n$；$u(x)$ 给出消费者从货篮 $x$ 中得到的效用的一个度量。对于两个货篮的选择，消费者将挑选具有较大效用的一个，当具有相同效用时，消费者心情矛盾。对 $u$ 是单调非减函数的假设是合理的，这表明对于任意商品，在其他商品数量一定的情况下，消费者总是倾向于获取更多的量。同样，假设 $u$ 是凹函数也是合理的。这对饱和进行了建模，或者说我们在增加商品的量的同时降低了边际效用。

现在假设给出了一些消费者偏好数据，但是我们不知道潜在的效用函数 $u$。具体地，我们有一个货篮集合 $a_1,\cdots,a_m\in[0,1]^n$ 和关于它们之间偏好度的信息：
\begin{equation}
u(a_i)>u(a_j)\ \text{对于 }(i,j)\in\mathcal{P},\qquad
u(a_i)\geq u(a_j)\ \text{对于 }(i,j)\in\mathcal{P}_{\mathrm{weak}},
\tag{6.22}
\end{equation}
其中 $\mathcal{P},\mathcal{P}_{\mathrm{weak}}\subseteq\{1,\cdots,m\}\times\{1,\cdots,m\}$ 是给定的。这里的 $\mathcal{P}$ 给出了已知偏好的集合：$(i,j)\in\mathcal{P}$ 表明已知货篮 $a_i$ 优于货篮 $a_j$。集合 $\mathcal{P}_{\mathrm{weak}}$ 给出了已知的弱偏好：$(i,j)\in\mathcal{P}_{\mathrm{weak}}$ 意味着货篮 $a_i$ 优于货篮 $a_j$，或者两者具有相同的吸引力。

首先考虑下面的问题：我们如何确定给定的数据是否相合，即是否存在凹的非减的效用函数使得 (6.22) 成立？这等价于求解可行性问题
\begin{equation}
\begin{aligned}
\text{find}\quad & u\\
\text{subject to}\quad & u:\mathbb{R}^n\to\mathbb{R}\text{ 凹和非减}\\
& u(a_i)>u(a_j),\quad (i,j)\in\mathcal{P}\\
& u(a_i)\geq u(a_j),\quad (i,j)\in\mathcal{P}_{\mathrm{weak}},
\end{aligned}
\tag{6.23}
\end{equation}
其中函数 $u$ 是（无穷维）优化变量。因为 (6.23) 中的约束是齐次的，所以，我们可以将问题表述为等价形式
\begin{equation}
\begin{aligned}
\text{find}\quad & u\\
\text{subject to}\quad & u:\mathbb{R}^n\to\mathbb{R}\text{ 凹和非减}\\
& u(a_i)\geq u(a_j)+1,\quad (i,j)\in\mathcal{P}\\
& u(a_i)\geq u(a_j),\quad (i,j)\in\mathcal{P}_{\mathrm{weak}},
\end{aligned}
\tag{6.24}
\end{equation}
这里只利用了不严格的不等式。（显然满足 (6.24) 的 $u$ 一定满足 (6.23)；反之，如果 $u$ 满足 (6.23)，那么可以通过伸缩变换使之满足 (6.24)。）此处利用第 326 页关于插值的结果，可以将该问题建模为（有限维）线性规划可行性问题：
\begin{equation}
\begin{aligned}
\text{find}\quad & u_1,\cdots,u_m,\quad g_1,\cdots,g_m\\
\text{subject to}\quad & g_i\succeq 0,\quad i=1,\cdots,m\\
& u_j\leq u_i+g_i^T(a_j-a_i),\quad i,j=1,\cdots,m\\
& u_i\geq u_j+1,\quad (i,j)\in\mathcal{P}\\
& u_i\geq u_j,\quad (i,j)\in\mathcal{P}_{\mathrm{weak}}.
\end{aligned}
\tag{6.25}
\end{equation}
通过求解这个线性规划可行性问题，我们可以确定是否存在凹的、非减的效用函数与给定的严格和不严格偏好集合相容。如果 (6.25) 可行，存在至少一个这样的效用函数（并且，实际上，我们可以从可行的 $u_1,\cdots,u_m$，$g_1,\cdots,g_m$ 构造出这样的分片线性函数）。如果 (6.25) 不可行，我们可以断定，没有凹的单调效用函数与给定的严格、不严格偏好集合相容。

作为一个例子，假设已知 $\mathcal{P}$ 和 $\mathcal{P}_{\mathrm{weak}}$ 是至少与一个凹的单增效用函数相合的偏好。考虑不在 $\mathcal{P}$ 或 $\mathcal{P}_{\mathrm{weak}}$ 中的 $(k,l)$，即消费者对于 $k$ 和 $l$ 的偏好未知。在某些情况下，我们甚至可以在不知道潜在的偏好函数的情况下可以断定货篮 $k$ 和 $l$ 的偏好。为此，我们用不等式 $u(a_k)\leq u(a_l)$ 来扩展已知偏好 (6.22)，这意味着 $l$ 优于 $k$，或者具有相同的吸引力。然后，我们在包含附加弱偏好 $u(a_k)\leq u(a_l)$ 的情况下求解线性规划可行性问题 (6.25)。如果扩展的偏好集合不可行，则意味着任意与给定的消费者偏好数据相容的凹的非减效用函数，都必须满足 $u(a_k)>u(a_l)$。换言之，我们可以在不知道潜在的效用函数的情况下，断定货篮 $k$ 优于货篮 $l$。

现在，我们利用消费者偏好数据（当然，不使用真实的效用函数 $u$）来将 $a_0=(0.5,0.5)$ 与这 40 个货篮分别进行比较。对于每一个原有的货篮 $a_i$，我们求解前述的线性规划可行性问题，用以检查我们是否可以判断货篮 $a_0$ 优于货篮 $a_i$。类似地，我们可以检查是否能得出 $a_i$ 优于货篮 $a_0$ 的结论。对于每个货篮 $a_i$，有三种可能的输出：我们能够确认 $a_0$ 一定优于 $a_i$，$a_i$ 一定优于 $a_0$，或者（如果两个 LP 可行性问题都可行）无法得到上述结论。（这里的一定优于意味着这个优越性对于任意与原始给定数据相容的凹的非减效用函数都成立。）

我们发现，货篮中的 21 个一定不如 $(0.5,0.5)$，而有 14 个一定优于它。对于剩余的 5 个货篮，我们无法从顾客偏好数据中得出任何结论。这些结果显示在图 6.26 中。注意，只利用效用函数的单调性，能够判断出在 $(0.5,0.5)$ 左下方的货篮一定不如 $(0.5,0.5)$；类似地，$(0.5,0.5)$ 右上方的点会更好。因此，对于这 17 个点，并没有必要去求解 LP 可行性问题 (6.25)。但是，对于处于另外两个区域象限的点，需要用到凹性假设并求解 LP 可行性问题 (6.25)。

\section*{参考文献}

Huber 在 [Hub64, Hub81] 中分析了利用不同罚函数进行逼近的鲁棒性，他同时提出了罚函数 (6.4)。对数障碍罚函数出现在控制领域，在系统的闭环频率响应中得到了应用，并且还有其他称呼，如中心 $H_\infty$ 或无风险控制；参见 Boyd 和 Barratt 的 [BB91] 和其中的参考文献。

很多书中都包含了正则化逼近的内容，如 Tikhonov 和 Arsenin 的 [TA77] 以及 Hansen 的 [Han98]。Tikhonov 正则化有时又称为岭回归（Golub 和 Van Loan 的 [GL89，第 564 页]）。$\ell_1$-范数正则化的最小二乘估计也称为套索（lasso）（Tibshirani [Tib96]）而被大家所知。其他最小二乘正则化和回归选择技术在 Hastie、Tibshirani 和 Friedman 的 [HTF01，\S3.4] 中进行了讨论和比较。

总变差去噪声方法是在图像重构中由 Rudin、Osher 和 Fatemi 在 [ROF92] 提出的。范数有界的不确定性下的鲁棒最小二乘问题（第 310 页）由 El Ghaoui 和 Lebret 在 [EL97] 以及 Chandrasekaran、Golub、Gu 和 Sayed 在 [CGGS98] 提出。El Ghaoui 和 Lebret 也给出了结构不确定下鲁棒最小二乘（第 312 页）的 SDP 公式。

Chen、Donoho 和 Saunders 在 [CDS01] 中讨论了利用线性规划的基筛选问题。他们称 $\ell_1$-范数正则化问题 (6.18) 为基筛选去除噪声。Meyer 和 Pratt 的 [MP68] 是早期的一篇关于效用函数界定问题的文章。

\section*{习题}

\subsection*{范数逼近与最小范数问题}

\noindent\textbf{6.1 对数障碍罚函数的二次界。}\quad
令 $\phi:\mathbb{R}\to\mathbb{R}$ 为具有极限 $a>0$ 的对数罚函数
\[
\phi(u)=
\begin{cases}
-a^2\log(1-(u/a)^2) & |u|<a\\
\infty & \text{其他情况}.
\end{cases}
\]
证明如果 $u\in\mathbb{R}^m$ 满足 $\|u\|_\infty<a$，那么
\[
\|u\|_2^2\leq \sum_{i=1}^m\phi(u_i)\leq
\frac{\phi(\|u\|_\infty)}{\|u\|_\infty^2}\|u\|_2^2.
\]
也就是说，如果 $\|u\|_\infty$ 相对于 $a$ 较小，那么 $\sum_{i=1}^m\phi(u_i)$ 能被 $\|u\|_2^2$ 很好地逼近。例如，如果 $\|u\|_\infty/a=0.25$，那么
\[
\|u\|_2^2\leq \sum_{i=1}^m\phi(u_i)\leq 1.033\cdot \|u\|_2^2.
\]

\noindent\textbf{6.2 常向量的 $\ell_1$-、$\ell_2$-和 $\ell_\infty$-范数逼近。}\quad
一个标量 $x\in\mathbb{R}$ 的范数逼近问题
\[
\begin{aligned}
\text{minimize}\quad & \|x\mathbf{1}-b\|
\end{aligned}
\]
对于 $\ell_1$-、$\ell_2$-和 $\ell_\infty$-范数的解是什么？

\noindent\textbf{6.3}\quad
将下面的问题构建为线性规划、二次规划、二阶锥规划或半定规划。问题数据为 $A\in\mathbb{R}^{m\times n}$ 和 $b\in\mathbb{R}^m$。$A$ 的行记为 $a_i^T$。

(a) 带有死区的罚函数逼近。$\text{minimize}\ \sum_{i=1}^m\phi(a_i^Tx-b_i)$，其中
\[
\phi(u)=
\begin{cases}
0 & |u|\leq a\\
|u|-a & |u|>a,
\end{cases}
\]
其中 $a>0$。

(b) 对数障碍罚函数逼近。$\text{minimize}\ \sum_{i=1}^m\phi(a_i^Tx-b_i)$，其中
\[
\phi(u)=
\begin{cases}
-a^2\log(1-(u/a)^2) & |u|<a\\
\infty & |u|\geq a,
\end{cases}
\]
其中 $a>0$。

(c) Huber 罚函数逼近。$\text{minimize}\ \sum_{i=1}^m\phi(a_i^Tx-b_i)$，其中
\[
\phi(u)=
\begin{cases}
u^2 & |u|\leq M\\
M(2|u|-M) & |u|>M,
\end{cases}
\]
其中 $M>0$。

(d) 对数-Chebyshev 逼近。$\text{minimize}\ \max_{i=1,\cdots,m}|\log(a_i^Tx)-\log b_i|$。我们假设 $b\succ 0$。等价的凸形式为
\[
\begin{aligned}
\text{minimize}\quad & t\\
\text{subject to}\quad & 1/t\leq a_i^Tx/b_i\leq t,\quad i=1,\cdots,m,
\end{aligned}
\]
其变量为 $x\in\mathbb{R}^n$ 和 $t\in\mathbb{R}$，定义域为 $\mathbb{R}^n\times\mathbb{R}_{++}$。

(e) 极小化最大的 $k$ 个残差之和：
\[
\begin{aligned}
\text{minimize}\quad & \sum_{i=1}^k |r|_{[i]}\\
\text{subject to}\quad & r=Ax-b,
\end{aligned}
\]
其中 $|r|_{[1]}\geq |r|_{[2]}\geq \cdots\geq |r|_{[m]}$ 为 $|r_1|,|r_2|,\cdots,|r_m|$ 的递减排序。（对于 $k=1$，这个问题退化为 $\ell_\infty$-范数逼近；对于 $k=m$，这个问题退化为 $\ell_1$-范数逼近。）提示：参见习题 5.19。

\noindent\textbf{6.4 $\ell_1$-范数逼近的可微近似。}\quad
带有参数 $\epsilon>0$ 的函数 $\phi(u)=(u^2+\epsilon)^{1/2}$ 有时被用来作为绝对值函数 $|u|$ 的可微近似。为了近似地求解 $\ell_1$-范数逼近问题
\begin{equation}
\begin{aligned}
\text{minimize}\quad & \|Ax-b\|_1,
\end{aligned}
\tag{6.26}
\end{equation}
其中 $A\in\mathbb{R}^{m\times n}$，我们求解替代问题
\begin{equation}
\begin{aligned}
\text{minimize}\quad & \sum_{i=1}^m \phi(a_i^Tx-b_i),
\end{aligned}
\tag{6.27}
\end{equation}
其中，$a_i^T$ 为 $A$ 的第 $i$ 行。假设 $\mathbf{rank}\ A=n$。

记 $\ell_1$-范数逼近问题 (6.26) 的最优值为 $p^\star$。用 $\hat{x}$ 代表近似问题 (6.27) 的最优解，并用 $\hat{r}$ 表示相应的残差 $\hat{r}=A\hat{x}-b$。

(a) 证明 $p^\star\geq \sum_{i=1}^m \hat{r}_i^2/(\hat{r}_i^2+\epsilon)^{1/2}$。

(b) 证明
\[
\|A\hat{x}-b\|_1\leq p^\star+\sum_{i=1}^m |\hat{r}_i|\left(1-\frac{|\hat{r}_i|}{(\hat{r}_i^2+\epsilon)^{1/2}}\right).
\]
（计算 $\hat{x}$ 后，通过评估右端，我们得到了 $\hat{x}$ 对于 $\ell_1$-范数逼近问题次优情况的界。）

\noindent\textbf{6.5 最小长度逼近。}\quad
考虑问题
\[
\begin{aligned}
\text{minimize}\quad & \mathbf{length}(x)\\
\text{subject to}\quad & \|Ax-b\|\leq \epsilon,
\end{aligned}
\]
其中 $\mathbf{length}(x)=\min\{k\mid x_i=0\text{ 对于 }i>k\}$。问题的变量为 $x\in\mathbb{R}^n$，参数为 $A\in\mathbb{R}^{m\times n}$、$b\in\mathbb{R}^m$ 和 $\epsilon>0$。在回归问题中，我们被要求给出以精度 $\epsilon$ 逼近向量 $b$ 所需要选用的 $A$ 的列的最少数量。

证明这是一个拟凸优化问题。

\noindent\textbf{6.6 一些罚函数逼近问题的对偶。}\quad
对于下列的罚函数 $\phi:\mathbb{R}\to\mathbb{R}$，导出问题
\[
\begin{aligned}
\text{minimize}\quad & \sum_{i=1}^m \phi(r_i)\\
\text{subject to}\quad & r=Ax-b
\end{aligned}
\]
的 Lagrange 对偶形式。这个问题的变量为 $x\in\mathbb{R}^n$ 和 $r\in\mathbb{R}^m$。

(a) 死区-线性罚函数（死区宽度 $a=1$），
\[
\phi(u)=
\begin{cases}
0 & |u|\leq 1\\
|u|-1 & |u|>1.
\end{cases}
\]

(b) Huber 罚函数（$M=1$），
\[
\phi(u)=
\begin{cases}
u^2 & |u|\leq 1\\
2|u|-1 & |u|>1.
\end{cases}
\]

(c) 对数障碍罚函数（极限为 $a=1$），
\[
\phi(u)=-\log(1-u^2),\qquad \mathbf{dom}\ \phi=(-1,1).
\]

(d) 距离 1 的相对偏差
\[
\phi(u)=\max\{u,1/u\}=
\begin{cases}
u & u\geq 1\\
1/u & u\leq 1,
\end{cases}
\]
其中 $\mathbf{dom}\ \phi=\mathbb{R}_{++}$。

\subsection*{正则化与鲁棒逼近}

\noindent\textbf{6.7 Euclid 范数下的双准则最优化。}\quad
我们考虑双准则优化问题
\[
\begin{aligned}
\text{minimize（关于 $\mathbb{R}_+^2$）}\quad & (\|Ax-b\|_2^2,\|x\|_2^2),
\end{aligned}
\]
其中 $A\in\mathbb{R}^{m\times n}$ 的秩为 $r$ 而 $b\in\mathbb{R}^m$。说明从 $A$ 的对角分解值
\[
A=U\operatorname{diag}(\sigma)V^T=\sum_{i=1}^r\sigma_i u_i v_i^T
\]
（参见 \S A.5.4），如何找到下面各个问题的解。

(a) Tikhonov 正则化：极小化 $\|Ax-b\|_2^2+\delta\|x\|_2^2$。

(b) 在 $\|x\|_2^2=\gamma$ 约束下极小化 $\|Ax-b\|_2^2$。

(c) 在 $\|x\|_2^2=\gamma$ 约束下极大化 $\|Ax-b\|_2^2$。

这里 $\delta$ 和 $\gamma$ 为正的参数。

你的结果提供了有效的算法用以计算最优权衡曲线和双准则问题的可达值集合。

\noindent\textbf{6.8}\quad
将下列鲁棒逼近问题建模为线性规划、二次规划、二阶锥规划或半定规划。对于每个子问题，分别考虑 $\ell_1$-、$\ell_2$-和 $\ell_\infty$-范数。

(a) 使用有限参数值集合的随机鲁棒逼近，即范数和问题
\[
\begin{aligned}
\text{minimize}\quad & \sum_{i=1}^k p_i\|A_ix-b\|,
\end{aligned}
\]
其中 $p\succeq 0$，$\mathbf{1}^Tp=1$。（参见 \S6.4.1。）

(b) 系数有界制的最坏情况鲁棒逼近：
\[
\begin{aligned}
\text{minimize}\quad & \sup_{A\in\mathcal{A}}\|Ax-b\|,
\end{aligned}
\]
其中
\[
\mathcal{A}=\{A\in\mathbb{R}^{m\times n}\mid l_{ij}\leq a_{ij}\leq u_{ij},\ i=1,\cdots,m,\ j=1,\cdots,n\}.
\]
这里对 $A$ 的分量给出上下界描述的不确定性集合。我们假设 $l_{ij}<u_{ij}$。

(c) 对于多面体不确定性的最坏情况鲁棒逼近：
\[
\begin{aligned}
\text{minimize}\quad & \sup_{A\in\mathcal{A}}\|Ax-b\|,
\end{aligned}
\]
其中
\[
\mathcal{A}=\{[a_1\ \cdots\ a_m]^T\mid C_i a_i\preceq d_i,\ i=1,\cdots,m\}.
\]
这里通过对每一行可能值给出的多面体 $\mathcal{P}_i=\{a_i\mid C_i a_i\preceq d_i\}$ 描述了不确定性。$C_i\in\mathbb{R}^{p_i\times n}$，$d_i\in\mathbb{R}^{p_i}$，$i=1,\cdots,m$ 为给定参数。我们假设多面体 $\mathcal{P}_i$ 是非空和有界的。

\subsection*{范数拟合与插值}

\noindent\textbf{6.9 极小极大有理函数拟合。}\quad
证明下面的问题是拟凸的：
\[
\begin{aligned}
\text{minimize}\quad & \max_{i=1,\cdots,k}\left|\frac{p(t_i)}{q(t_i)}-y_i\right|,
\end{aligned}
\]
其中
\[
p(t)=a_0+a_1t+a_2t^2+\cdots+a_mt^m,\qquad
q(t)=1+b_1t+\cdots+b_nt^n,
\]
并且目标函数的定义域为
\[
D=\{(a,b)\in\mathbb{R}^{m+1}\times\mathbb{R}^n\mid q(t)>0,\ \alpha\leq t\leq\beta\}.
\]
在这个问题中，我们用一个有理函数 $p(t)/q(t)$ 来拟合给定数据，同时约束分母多项式在区间 $[\alpha,\beta]$ 上为正。优化变量为分子和分母的系数 $a_i$、$b_i$。插值点 $t_i\in[\alpha,\beta]$ 和需要的函数值 $y_i$，$i=1,\cdots,k$ 给定。

\noindent\textbf{6.10 利用凹的非负非减二次函数拟合数据。}\quad
给定数据
\[
x_1,\cdots,x_N\in\mathbb{R}^n,\qquad y_1,\cdots,y_N\in\mathbb{R},
\]
我们希望拟合得到具有下列形式的二次函数
\[
f(x)=(1/2)x^TPx+q^Tx+r,
\]
其中，$P\in\mathbf{S}^n$、$q\in\mathbb{R}^n$ 和 $r\in\mathbb{R}$ 为模型参数（因此，是拟合问题中的变量）。

我们的模型将仅仅被用在区域 $B=\{x\in\mathbb{R}^n\mid l\preceq x\preceq u\}$ 中。你可以假设 $l\prec u$，给定数据 $x_i$ 处在这个区域。我们用简单的平方和误差目标
\[
\sum_{i=1}^N(f(x_i)-y_i)^2
\]
作为拟合的准则。我们向函数 $f$ 添加一些约束。首先，它必须是凹的。其次，它必须在 $B$ 中非负，即对于所有的 $z\in B$ 均有 $f(z)\geq 0$。第三，$f$ 必须是 $B$ 上的非减函数，即只要 $z,\bar{z}\in B$ 满足 $z\succeq \bar{z}$，我们都有 $f(z)\leq f(\bar{z})$。

说明如何将这个拟合问题建模为凸问题。尽可能简化你的式子。

\noindent\textbf{6.11 最小二乘方向插值。}\quad
设 $F_1,\cdots,F_n:\mathbb{R}^k\to\mathbb{R}^p$，我们构造线性组合 $F:\mathbb{R}^k\to\mathbb{R}^p$，
\[
F(u)=x_1F_1(u)+\cdots+x_nF_n(u),
\]
其中，$x$ 为插值问题中的变量。

在这个问题中，我们要求 $\angle(F(v_j),q_j)=0$，$j=1,\cdots,m$，其中 $q_j$ 为 $\mathbb{R}^p$ 中的给定向量，并且我们假设其满足 $\|q_j\|_2=1$。换言之，我们要求 $F$ 的方向在点 $v_j$ 上取特定的值。为保证 $F(v_j)$ 是非零的（值为零将使得角度无定义），我们对最小长度加以约束 $\|F(v_j)\|_2\geq \epsilon$，$j=1,\cdots,m$，其中 $\epsilon>0$ 为给定数据。

说明如何使用凸优化，找到 $x$ 以极小化 $\|x\|^2$ 并且满足上述方向（和最小长度）条件。

\noindent\textbf{6.12 单调函数插值。}\quad
函数 $f:\mathbb{R}^k\to\mathbb{R}$（关于 $\mathbb{R}_+^k$）是单调非减的，如果只要 $u\succeq v$ 就有 $f(u)\geq f(v)$。

(a) 证明存在单调非减函数 $f:\mathbb{R}^k\to\mathbb{R}$ 对于 $i=1,\cdots,m$ 满足 $f(u_i)=y_i$ 的充要条件是，
\[
y_i\geq y_j\quad \text{只要 }u_i\succeq u_j,\quad i,j=1,\cdots,m.
\]

(b) 证明当且仅当存在 $g_i\in\mathbb{R}^k$，$i=1,\cdots,m$ 使得
\[
g_i\succeq 0,\quad i=1,\cdots,m,\qquad
y_j\geq y_i+g_i^T(u_j-u_i),\quad i,j=1,\cdots,m
\]
时，存在凸的单调非减函数 $f:\mathbb{R}^k\to\mathbb{R}$，其中 $\mathbf{dom}\ f=\mathbb{R}^k$，满足对于 $i=1,\cdots,m$ 有 $f(u_i)=y_i$。

\noindent\textbf{6.13 拟凸函数插值。}\quad
证明存在拟凸函数 $f:\mathbb{R}^k\to\mathbb{R}$ 对 $i=1,\cdots,m$ 满足 $f(u_i)=y_i$ 的充要条件是，存在 $g_i\in\mathbb{R}^k$，$i=1,\cdots,m$ 使得
\[
g_i^T(u_j-u_i)\leq -1\quad \text{只要 } y_j<y_i,\quad i,j=1,\cdots,m.
\]

\noindent\textbf{6.14 [Nes00] 正实部函数插值。}\quad
设 $z_1,\cdots,z_n\in\mathbb{C}$ 为 $n$ 个不同的点，并且 $|z_i|>1$。我们定义 $K_{\mathrm{np}}$ 为向量 $y\in\mathbb{C}^n$ 的集合，对于其中每个向量，存在函数 $f:\mathbb{C}\to\mathbb{C}$ 满足下面的条件：
\begin{itemize}
\item $f$ 是正实部函数，这意味着它在单位圆外部（即 $|z|>1$）是解析的，并且在单位圆外实部非负（对于 $|z|>1$，有 $\Re f(z)\geq 0$）。
\item $f$ 满足插值条件
\[
f(z_1)=y_1,\qquad f(z_2)=y_2,\qquad \cdots,\qquad f(z_n)=y_n.
\]
\end{itemize}
如果我们将正实部函数记为 $\mathcal{F}$，那么 $K_{\mathrm{np}}$ 可以表示为
\[
K_{\mathrm{np}}=\{y\in\mathbb{C}^n\mid \exists f\in\mathcal{F},\ y_k=f(z_k),\ k=1,\cdots,n\}.
\]

(a) 可以证明 $f$ 是正实部的充要条件是，存在一个非减函数 $\rho$ 使得对于所有满足 $|z|>1$ 的 $z$，都有
\[
f(z)=i\Im f(\infty)+\int_0^{2\pi}\frac{e^{i\theta}+z^{-1}}{e^{i\theta}-z^{-1}}\,d\rho(\theta),
\]
其中 $i=\sqrt{-1}$（参见 [KN77，第 389 页]）。用这个表达式来证明 $K_{\mathrm{np}}$ 是一个闭凸锥。

(b) 我们将利用 $\Re(x^Hy)$ 和 $x,y\in\mathbb{C}^n$ 之间的内积，其中 $x^H$ 表示 $x$ 的复共轭转置。证明 $K_{\mathrm{np}}$ 的对偶锥由
\[
K_{\mathrm{np}}^*=
\left\{x\in\mathbb{C}^n\ \middle|\ 
\Im(\mathbf{1}^Tx)=0,\ 
\Re\left(\sum_{l=1}^n x_l\frac{e^{-i\theta}+\bar{z}_l^{-1}}{e^{-i\theta}-\bar{z}_l^{-1}}\right)\geq 0,\ \forall\theta\in[0,2\pi]
\right\}
\]
给出。

(c) 证明
\[
K_{\mathrm{np}}^*=
\left\{x\in\mathbb{C}^n\ \middle|\ 
\exists Q\in\mathbf{H}_+^n,\ x_l=\sum_{k=1}^n\frac{Q_{kl}}{1-z_k^{-1}\bar{z}_l^{-1}},\ l=1,\cdots,n
\right\},
\]
其中 $\mathbf{H}_+^n$ 表示 $n\times n$ 的半正定 Hermitian 矩阵。

利用下面的结果（即 Riesz--Fejer 定理，参见 [KN77，第 60 页]）。形如
\[
\sum_{k=0}^n(y_ke^{-ik\theta}+\bar{y}_ke^{ik\theta})
\]
的函数对于所有 $\theta$ 非负的充要条件是，存在 $a_0,\cdots,a_n\in\mathbb{C}$ 使得
\[
\sum_{k=0}^n(y_ke^{-ik\theta}+\bar{y}_ke^{ik\theta})
=\left|\sum_{k=0}^n a_ke^{ik\theta}\right|^2.
\]

(d) 证明 $K_{\mathrm{np}}=\{y\in\mathbb{C}^n\mid P(y)\succeq 0\}$，其中 $P(y)\in\mathbf{H}^n$ 定义为
\[
P(y)_{kl}=\frac{y_k+\bar{y}_l}{1-z_k^{-1}\bar{z}_l^{-1}},\qquad l,k=1,\cdots,n.
\]
矩阵 $P(y)$ 被称为关于点 $z_k$，$y_k$ 的 Nevanlinna-Pick 矩阵。

提示：如果我们在 (a) 中所述，$K_{\mathrm{np}}$ 是一个闭凸锥，因此 $K_{\mathrm{np}}=K_{\mathrm{np}}^{**}$。

(e) 作为应用，将下面问题建模为一个凸优化问题：
\[
\begin{aligned}
\text{minimize}\quad & \sum_{k=1}^n |f(z_k)-w_k|^2\\
\text{subject to}\quad & f\in\mathcal{F}.
\end{aligned}
\]
问题的数据为 $n$ 个点 $z_k$ 和 $n$ 个复数 $w_1,\cdots,w_n$，并且 $|z_k|>1$。我们在所有正实部函数中优化 $f$。

\end{document}
